\documentclass[11pt,a4paper]{article}
\usepackage[utf8]{inputenc}
% lmodern must precede fontenc: without a scalable Type1 font the T1 encoding
% falls back to bitmap fonts, and microtype's expansion then aborts the run
% with "auto expansion is only possible with scalable fonts".
\usepackage{lmodern}
\usepackage[T1]{fontenc}
\usepackage[margin=1in]{geometry}
\usepackage{amsmath,amssymb}
\usepackage{graphicx}
\usepackage{microtype}
\usepackage[hidelinks,breaklinks]{hyperref}
\usepackage{url}
\urlstyle{same}

% Generated by code/tools/md_to_tex.py from the SurveyForge Markdown output.
% Citation numbers match the source .md exactly.

\title{Comprehensive Survey of Visual Adversarial Attacks and Defenses in the Physical World}
\author{Generated by SurveyForge}
\date{}

\begin{document}
\maketitle
\tableofcontents
\newpage



\section{Introduction}


The concept of visual adversarial attacks has garnered significant attention in recent years, particularly in the context of physical world applications. These attacks involve manipulating real-world objects or environments to deceive computer vision systems, which can have severe consequences for security, safety, and reliability \cite{doi:10.1201/9781351251389-8}. For instance, an attacker can create a physical adversarial example by printing a specially designed sticker and attaching it to a stop sign, causing a self-driving car to misclassify the sign and potentially leading to accidents \cite{arxiv:1707.08945}. The history of visual adversarial attacks dates back to the discovery of digital adversarial examples, which are small perturbations added to digital images to fool image classifiers \cite{doi:10.24963/ijcai.2021/694}. However, the extension of these attacks to the physical world has raised new challenges and concerns, as physical adversarial examples can be more difficult to detect and defend against \cite{doi:10.1109/cvpr42600.2020.01426}.

The significance of visual adversarial attacks in physical world applications cannot be overstated. Autonomous vehicles, surveillance systems, and facial recognition systems are just a few examples of computer vision systems that can be compromised by these attacks \cite{doi:10.1016/j.sysarc.2020.101766}. The potential consequences of such attacks are dire, ranging from financial loss to loss of human life \cite{doi:10.1109/jiot.2021.3099164}. Therefore, it is essential to understand and defend against these attacks to ensure the security, safety, and reliability of computer vision systems \cite{doi:10.1016/j.eng.2019.12.012}. Researchers have proposed various defense mechanisms, including adversarial training, input validation, and detection-based approaches \cite{arxiv:2102.01356}. However, the development of effective defense mechanisms is an ongoing challenge, as attackers can adapt and evolve their methods to bypass existing defenses \cite{doi:10.1109/iros51168.2021.9636638}.

Recent studies have demonstrated the effectiveness of physical adversarial attacks in various scenarios, including object detection \cite{arxiv:1906.11897} and facial recognition \cite{doi:10.1609/aaai.v32i1.12341}. These attacks can be launched using various methods, including printing or projecting adversarial patterns onto objects or environments \cite{arxiv:2007.04137}. The use of deep learning models in computer vision systems has increased the vulnerability of these systems to adversarial attacks \cite{doi:10.1109/access.2022.3208131}. Therefore, it is crucial to develop robust defense mechanisms that can detect and mitigate these attacks \cite{doi:10.1109/asp-dac47756.2020.9045584}.

The analysis of visual adversarial attacks in physical world applications requires a comprehensive understanding of the underlying concepts, including the types of attacks, the methods used to launch these attacks, and the defense mechanisms available to mitigate them \cite{doi:10.1007/s11633-019-1211-x}. Researchers have proposed various taxonomies for categorizing visual adversarial attacks, including digital and physical attacks \cite{doi:10.24963/ijcai.2021/694}. Digital attacks involve adding perturbations to digital images, while physical attacks involve manipulating real-world objects or environments \cite{arxiv:1707.08945}. The development of effective defense mechanisms requires a deep understanding of the strengths and limitations of each approach \cite{arxiv:2102.01356}.

In conclusion, visual adversarial attacks in physical world applications pose a significant threat to the security, safety, and reliability of computer vision systems. The development of effective defense mechanisms requires a comprehensive understanding of the underlying concepts, including the types of attacks, the methods used to launch these attacks, and the defense mechanisms available to mitigate them \cite{doi:10.1016/j.eng.2019.12.012}. Researchers must continue to investigate and develop new defense mechanisms to stay ahead of the evolving threat landscape \cite{doi:10.1145/3459637.3482029}. By synthesizing information from various studies, including \cite{doi:10.24963/ijcai.2021/694}, \cite{arxiv:1707.08945}, and \cite{doi:10.1109/access.2022.3208131}, we can gain a deeper understanding of the challenges and opportunities in this field, ultimately contributing to the development of more robust and secure computer vision systems.


\section{Fundamentals of Adversarial Attacks}



\subsection{Introduction to Physical Adversarial Attacks}


Physical adversarial attacks represent a significant threat to the security and reliability of computer vision systems in the physical world. These attacks involve manipulating real-world objects or environments to deceive computer vision systems, which can have severe consequences in safety-critical applications such as autonomous vehicles, surveillance systems, and facial recognition systems. As highlighted in \cite{arxiv:1707.07397}, standard methods for generating adversarial examples do not consistently fool neural network classifiers in the physical world due to various transformations and perturbations. However, researchers have demonstrated the existence of robust 3D adversarial objects that can deceive computer vision systems in the physical world.

The concept of physical adversarial attacks is closely related to the idea of adversarial examples in the digital domain. Adversarial examples are inputs to a machine learning model that are specifically designed to cause the model to make a mistake. In the digital domain, adversarial examples can be generated using various techniques such as gradient-based methods or evolutionary algorithms. However, generating physical adversarial attacks is more challenging due to the need to consider real-world constraints such as lighting, viewpoint, and environmental conditions. As discussed in \cite{arxiv:1707.08945}, physical adversarial attacks can be generated using techniques such as printing or projecting, and can be designed to be robust to various transformations and perturbations.

One of the key challenges in generating physical adversarial attacks is ensuring that the attack is robust to various environmental conditions. As noted in \cite{doi:10.24963/ijcai.2021/694}, adversarial examples can be sensitive to transformations such as rotation, scaling, and lighting changes. To address this challenge, researchers have proposed various techniques such as Expectation Over Transformation (EOT) to improve the robustness of physical adversarial attacks. EOT involves generating adversarial examples that are robust to various transformations by optimizing over a distribution of transformations.

Physical adversarial attacks can be categorized into different types based on the level of access to the target model and the type of attack. As discussed in \cite{doi:10.1109/access.2018.2807385}, physical adversarial attacks can be classified into white-box, black-box, and gray-box attacks. White-box attacks assume that the attacker has full access to the target model, while black-box attacks assume that the attacker has no access to the target model. Gray-box attacks assume that the attacker has some knowledge about the target model but not full access.

The evaluation of physical adversarial attacks is crucial to understanding their effectiveness and robustness. As highlighted in \cite{doi:10.1109/jiot.2021.3099164}, the evaluation of physical adversarial attacks should consider various metrics such as attack success rate, robustness to transformations, and perceptual quality. The evaluation should also consider various scenarios and environmental conditions to ensure that the attack is robust and effective in different situations.

In recent years, various defense mechanisms have been proposed to mitigate the effects of physical adversarial attacks. As discussed in \cite{doi:10.1016/j.eng.2019.12.012}, defense mechanisms such as adversarial training, input validation, and detection-based approaches can be effective in improving the robustness of computer vision systems to physical adversarial attacks. However, the development of effective defense mechanisms is an ongoing challenge, and more research is needed to address the evolving threat of physical adversarial attacks.

In conclusion, physical adversarial attacks pose a significant threat to the security and reliability of computer vision systems in the physical world. Understanding the concept of physical adversarial attacks and their generation techniques is crucial to developing effective defense mechanisms. As highlighted in \cite{doi:10.1609/aaai.v32i1.12341}, the development of robust defense mechanisms requires a comprehensive understanding of the strengths and limitations of various approaches. Future research directions should focus on developing more effective defense mechanisms and evaluating their robustness to various physical adversarial attacks. As noted in \cite{doi:10.1109/cvpr42600.2020.00108}, the development of effective defense mechanisms is critical to ensuring the security and reliability of computer vision systems in safety-critical applications.


\subsection{Methods for Generating Physical Adversarial Examples}


Generating physical adversarial examples is a complex task that requires careful consideration of various factors, including the type of attack, the target model, and the environment in which the attack will be deployed. As \cite{arxiv:1707.07397} have shown, standard methods for generating adversarial examples do not consistently fool neural network classifiers in the physical world due to a combination of viewpoint shifts, camera noise, and other natural transformations. This challenge is closely related to the concept of physical adversarial attacks, which was discussed in the previous section, highlighting the need for robust and effective methods to generate physical adversarial examples.

To address this challenge, researchers have proposed various methods for generating physical adversarial examples, including printing, projecting, and other techniques. One approach to generating physical adversarial examples is to use printing-based methods. For example, \cite{arxiv:1707.08945} have demonstrated the existence of robust 3D adversarial objects that can fool image classifiers in the physical world. They propose a method for generating physical adversarial examples by printing adversarial images and then taking pictures of them. This approach has been shown to be effective in fooling image classifiers in various environments, including indoor and outdoor settings.

In addition to printing-based methods, projection-based methods have also been proposed for generating physical adversarial examples. For example, \cite{arxiv:1904.00759} have proposed a method for generating physical adversarial examples by projecting adversarial images onto objects in the physical world. This approach has been shown to be effective in fooling object detectors and other computer vision models. Other techniques, such as creating adversarial objects that can be placed in the physical world, have also been proposed, as seen in \cite{doi:10.1007/978-3-030-10925-7_4}.

The effectiveness of physical adversarial examples in different scenarios is a topic of ongoing research. For example, \cite{arxiv:1906.11897} have demonstrated that physical adversarial patches can be used to fool object detectors in various environments, including indoor and outdoor settings. However, \cite{arxiv:1707.03501} have argued that physical adversarial examples may not be a significant threat to autonomous vehicles, as they can be easily detected and mitigated using various defense mechanisms. These findings highlight the importance of considering the specific use case and environment when evaluating the effectiveness of physical adversarial examples.

Despite the potential effectiveness of physical adversarial examples, there are several challenges and limitations associated with generating and deploying them. For example, \cite{doi:10.1109/lcomm.2019.2901469} have demonstrated that physical adversarial examples can be used to attack end-to-end autoencoder communication systems, but the effectiveness of such attacks depends on various factors, including the type of attack and the environment in which it is deployed. These challenges and limitations will be further discussed in the following section, which will explore the robustness of physical adversarial examples to various transformations and perturbations.

In conclusion, generating physical adversarial examples is a complex task that requires careful consideration of various factors, including the type of attack, the target model, and the environment in which the attack will be deployed. Various methods have been proposed for generating physical adversarial examples, including printing, projecting, and other techniques. While physical adversarial examples have been shown to be effective in fooling image classifiers and object detectors in various environments, there are several challenges and limitations associated with generating and deploying them. As \cite{doi:10.1016/j.eng.2019.12.012} have noted, the development of effective defense mechanisms against physical adversarial examples is an ongoing area of research, and further studies are needed to fully understand the potential risks and consequences of such attacks. Future research directions may include the development of more robust and effective methods for generating physical adversarial examples, as well as the investigation of new defense mechanisms and countermeasures against such attacks, as discussed in \cite{doi:10.1049/cit2.12028} and \cite{doi:10.1109/access.2022.3208131}.


\subsection{Factors Influencing the Success of Physical Adversarial Attacks}


The success of physical adversarial attacks is influenced by a multitude of factors, including environmental conditions, object geometry, and sensor characteristics. Environmental conditions, such as lighting, viewpoint, and occlusions, play a significant role in determining the effectiveness of physical adversarial attacks \cite{doi:10.24963/ijcai.2021/694}. For instance, attacks that are successful in a controlled laboratory setting may not be as effective in outdoor environments with varying lighting conditions \cite{arxiv:1707.08945}. Furthermore, the geometry of the object being attacked can also impact the success of the attack. For example, objects with complex textures or shapes may be more difficult to attack than those with simpler geometries \cite{doi:10.1007/978-3-030-10925-7_4}.

Sensor characteristics, such as camera resolution and sensor sensitivity, can also influence the success of physical adversarial attacks \cite{arxiv:1906.11897}. For example, high-resolution cameras may be more susceptible to attacks that exploit small perturbations in the input image \cite{doi:10.1109/iccvw.2019.00257}. Additionally, the type of sensor used can also impact the effectiveness of the attack. For instance, attacks that are successful against cameras with high sensor sensitivity may not be as effective against cameras with lower sensitivity \cite{arxiv:1807.07769}.

The robustness of physical adversarial examples to various transformations and perturbations is also an important factor in determining the success of physical adversarial attacks \cite{doi:10.1609/aaai.v34i01.5443}. For example, attacks that are robust to rotations, scaling, and lighting changes may be more effective in real-world scenarios than those that are not \cite{arxiv:1802.06430}. Furthermore, the use of expectation over transformation (EOT) techniques can improve the robustness of physical adversarial examples to various transformations and perturbations \cite{doi:10.1109/cvpr42600.2020.01426}.

In addition to these factors, the success of physical adversarial attacks can also be influenced by the type of attack being used. For example, patch-based attacks, which involve adding a small patch to the object being attacked, may be more effective than object-based attacks, which involve modifying the entire object \cite{doi:10.1109/biosig52210.2021.9548291}. Additionally, the use of ensemble methods, which involve combining multiple models and attacks, can improve the effectiveness of physical adversarial attacks \cite{doi:10.1109/asp-dac47756.2020.9045584}.

The evaluation of physical adversarial attacks is also an important consideration. Traditional evaluation metrics, such as accuracy and robustness, may not be sufficient to fully capture the effectiveness of physical adversarial attacks \cite{doi:10.1007/978-3-030-58589-1_3}. For example, metrics that take into account the physical properties of the attack, such as the size and shape of the perturbation, may be more effective in evaluating the success of physical adversarial attacks \cite{arxiv:2007.04137}.

In conclusion, the success of physical adversarial attacks is influenced by a complex interplay of factors, including environmental conditions, object geometry, sensor characteristics, and the type of attack being used. By understanding these factors and developing more effective evaluation metrics, researchers can improve the robustness of deep learning models to physical adversarial attacks and develop more effective defense strategies \cite{arxiv:2001.05574}. Future research should focus on developing more robust and effective physical adversarial attacks, as well as improving the evaluation metrics used to assess their effectiveness \cite{doi:10.1109/tip.2021.3092822}. Additionally, the development of more effective defense strategies, such as adversarial training and input validation, can help to mitigate the risks associated with physical adversarial attacks \cite{doi:10.1109/cvpr42600.2020.00049}.


\subsection{Robustness of Physical Adversarial Examples to Transformations and Perturbations}


The robustness of physical adversarial examples to various transformations and perturbations is a critical aspect of adversarial attacks in the physical world. As \cite{arxiv:1707.07397} and \cite{arxiv:1707.08945} have shown, creating robust physical adversarial examples that can withstand different environmental conditions, such as weather, noise, and viewpoint changes, is essential for successful attacks. This is because physical adversarial examples are sensitive to various transformations and perturbations, which can affect their effectiveness. Therefore, improving the robustness of physical adversarial examples is crucial for ensuring the success of physical adversarial attacks.

To address this challenge, several approaches have been proposed. One approach is to use techniques such as Expectation Over Transformation (EOT) \cite{arxiv:1707.08945}, which involves simulating different transformations and perturbations during the generation process. This approach has been shown to improve the robustness of physical adversarial examples to various environmental conditions. Another approach is to use 3D modeling and simulation to craft adversarial patches that can withstand different transformations and perturbations \cite{doi:10.1016/j.neucom.2022.05.031}. These approaches have been used in various studies, including \cite{doi:10.1007/978-3-030-10925-7_4} and \cite{arxiv:1906.11897}, to generate robust physical adversarial examples.

A comparative analysis of different approaches reveals that each has its strengths and limitations. For example, \cite{doi:10.1007/978-3-030-10925-7_4} uses a combination of EOT and 3D modeling to generate robust physical adversarial examples, while \cite{arxiv:1906.11897} uses a patch-based approach to generate adversarial examples that can withstand different transformations and perturbations. However, as \cite{arxiv:1707.03501} notes, the effectiveness of physical adversarial examples can be limited by various factors, such as the complexity of the scene and the quality of the camera. Evaluating the strengths, limitations, and trade-offs of different approaches is crucial for understanding the robustness of physical adversarial examples.

Recent studies have also explored emerging trends and challenges in this area. For instance, \cite{arxiv:1910.11099} shows that using a T-shirt as an adversarial example can be effective in evading person detectors, but may not be robust to different transformations and perturbations. In contrast, \cite{doi:10.1109/cvpr42600.2020.01426} uses a generative adversarial network (GAN) to generate physical adversarial examples that can withstand different environmental conditions. Additionally, \cite{doi:10.1109/cvpr52688.2022.01487} and \cite{arxiv:2007.04137} propose new techniques for generating robust physical adversarial examples. Furthermore, \cite{doi:10.1109/cvpr46437.2021.00614} notes that backdoor attacks in the physical world pose a significant threat to the security of deep learning systems, and require the development of more effective defense mechanisms.

In conclusion, the robustness of physical adversarial examples to various transformations and perturbations is a critical aspect of adversarial attacks in the physical world. By synthesizing information from various studies, including \cite{doi:10.1201/9781351251389-8}, \cite{doi:10.1145/3319535.3339815}, and \cite{arxiv:2007.16118}, we can gain a deeper understanding of the challenges and opportunities in this area. Future research directions include the development of more sophisticated techniques for generating robust physical adversarial examples, as well as the creation of more effective defense mechanisms against backdoor attacks in the physical world. As \cite{doi:10.1109/asp-dac47756.2020.9045584} notes, developing robust defense mechanisms against physical adversarial attacks is crucial for ensuring the security and reliability of deep learning systems in the physical world.


\subsection{Attack Vectors and Threat Models for Physical Adversarial Attacks}


Physical adversarial attacks pose a significant threat to the security and reliability of computer vision systems in the physical world. These attacks involve manipulating real-world objects or environments to deceive computer vision models, and can be launched using various attack vectors and threat models. In this subsection, we will explore the different attack vectors and threat models used in physical adversarial attacks, including white-box, black-box, and gray-box attacks.

White-box attacks assume that the attacker has full knowledge of the target model's architecture, weights, and training data \cite{doi:10.1145/3128572.3140444}. This allows the attacker to craft highly effective adversarial examples that can mislead the model with high confidence. However, white-box attacks are less realistic in practice, as attackers often do not have access to the target model's internal workings. In contrast, black-box attacks assume that the attacker has no knowledge of the target model's architecture or weights, and must rely on indirect methods to craft adversarial examples \cite{arxiv:1904.02884}. Black-box attacks are more realistic, but often less effective than white-box attacks.

Gray-box attacks fall somewhere in between, assuming that the attacker has some knowledge of the target model's architecture or weights, but not full access \cite{doi:10.1109/cvpr42600.2020.00108}. Gray-box attacks can be more effective than black-box attacks, but less effective than white-box attacks. Regardless of the attack vector, physical adversarial attacks can be launched using various threat models, including patch-based, object-based, and scene-based attacks \cite{arxiv:1707.08945}. Patch-based attacks involve adding a small perturbation to a specific region of the input image, while object-based attacks involve modifying the appearance of a specific object in the scene. Scene-based attacks involve modifying the entire scene, such as by changing the lighting or weather conditions.

Recent studies have demonstrated the effectiveness of physical adversarial attacks in various scenarios, including autonomous driving \cite{doi:10.1109/cvpr42600.2020.01426}, face recognition \cite{doi:10.1016/j.patcog.2022.109009}, and object detection \cite{doi:10.1007/978-3-030-10925-7_4}. These attacks have been shown to be highly effective, even in the presence of various defenses and countermeasures \cite{doi:10.1109/cvpr46437.2021.00614}. To defend against physical adversarial attacks, researchers have proposed various methods, including adversarial training \cite{arxiv:1904.13000}, input validation \cite{doi:10.1609/aaai.v32i1.11828}, and detection-based approaches \cite{doi:10.1109/tnnls.2021.3105238}.

Despite these efforts, physical adversarial attacks remain a significant threat to the security and reliability of computer vision systems. As computer vision systems become increasingly ubiquitous in the physical world, the risk of physical adversarial attacks will only continue to grow. Therefore, it is essential to continue researching and developing effective defenses against these attacks. One potential direction for future research is the development of more robust and generalizable defenses that can withstand various types of physical adversarial attacks \cite{arxiv:1902.06705}. Another direction is the exploration of new attack vectors and threat models, such as the use of natural phenomena like shadows and lighting to launch physical adversarial attacks \cite{arxiv:2203.03818}.

In conclusion, physical adversarial attacks pose a significant threat to the security and reliability of computer vision systems in the physical world. Understanding the various attack vectors and threat models used in these attacks is essential for developing effective defenses. By continuing to research and develop new defenses, we can help mitigate the risk of physical adversarial attacks and ensure the safe and reliable operation of computer vision systems in the physical world \cite{arxiv:2209.14262}. As noted in \cite{doi:10.24963/ijcai.2021/694}, the study of physical adversarial attacks is crucial for promoting the application of AI techniques and ensuring the security development of AI. Furthermore, \cite{doi:10.1016/j.jisa.2020.102694} highlights the importance of generating robust and natural physical adversarial examples to evaluate the robustness of object detectors. Overall, the development of effective defenses against physical adversarial attacks requires a deep understanding of the various attack vectors and threat models used in these attacks, as well as the development of new and innovative defense strategies \cite{doi:10.1109/asp-dac47756.2020.9045584}.


\subsection{Evaluation Metrics and Benchmarking for Physical Adversarial Attacks}


The development of robust evaluation metrics and benchmarking protocols is crucial for assessing the effectiveness of physical adversarial attacks, as it enables the comparison of different attack methods and defense strategies. As highlighted in \cite{arxiv:1712.07107}, the lack of standardized evaluation metrics hinders the comparison of different attack methods and defense strategies, which is essential for advancing the field of physical adversarial attacks. To address this, researchers have proposed various evaluation metrics, including success rate, distance metric (e.g., L2, L$\infty$), and perceptual path length \cite{arxiv:1712.09665}, which aim to capture the effectiveness, stealthiness, and robustness of physical adversarial attacks.

The choice of evaluation metric depends on the specific attack scenario and the characteristics of the target model \cite{doi:10.1109/access.2018.2807385}. For instance, in the context of object detection, metrics such as average precision (AP) and mean average precision (mAP) are commonly used \cite{doi:10.1109/lcomm.2019.2901469}, as they provide a comprehensive understanding of the attack's impact on the model's performance. In contrast, for image classification tasks, metrics such as top-1 error rate and top-5 error rate are more relevant \cite{arxiv:1911.05268}, as they capture the model's ability to correctly classify images.

In addition to evaluation metrics, benchmarking protocols are also essential for evaluating the effectiveness of physical adversarial attacks. \cite{doi:10.1007/978-3-030-58589-1_3} presents a comprehensive dataset of physical adversarial attacks on object detection models, which can be used to evaluate the robustness of different models and defense strategies. Similarly, \cite{doi:10.1109/tgrs.2022.3225306} proposes a benchmarking framework for evaluating the effectiveness of adversarial patches against aerial detection models, highlighting the importance of standardized benchmarking protocols.

Despite the progress made in developing evaluation metrics and benchmarking protocols, there are still several challenges and limitations that need to be addressed. For example, \cite{doi:10.1007/s11633-019-1211-x} highlights the need for more realistic and comprehensive evaluation metrics that capture the complexity of physical-world scenarios, such as the impact of environmental factors on the attack's effectiveness. Additionally, \cite{doi:10.1109/cvpr42600.2020.00108} emphasizes the importance of considering the trade-offs between attack effectiveness and stealthiness, as highly effective attacks may not be stealthy, and vice versa.

To overcome these challenges, researchers have proposed various approaches, including the use of transferability attacks \cite{arxiv:2104.02361} and ensemble methods \cite{doi:10.1109/asp-dac47756.2020.9045584}, which can improve the robustness of defense strategies against physical adversarial attacks. Moreover, \cite{doi:10.1109/cvpr52688.2022.01487} presents a novel framework for generating physical adversarial examples that can camouflage objects in the physical world, highlighting the importance of developing more sophisticated attack methods to test the robustness of defense strategies.

In conclusion, the development of robust evaluation metrics and benchmarking protocols is essential for assessing the effectiveness of physical adversarial attacks, and for advancing the field of physical adversarial attacks. While significant progress has been made in this area, there are still several challenges and limitations that need to be addressed, such as the development of more realistic and comprehensive evaluation metrics, and the improvement of the robustness of defense strategies against physical adversarial attacks. By addressing these challenges, we can improve the robustness of machine learning models and mitigate the risks associated with physical adversarial attacks \cite{arxiv:2209.09502}, ultimately ensuring the security and reliability of machine learning models in the physical world \cite{doi:10.1016/j.patcog.2022.109009}.


\section{Threat Models and Attack Vectors}



\subsection{Introduction to Threat Models}


Threat models play a crucial role in understanding and defending against visual adversarial attacks, which pose significant risks to the security, safety, and reliability of computer vision systems \cite{doi:10.1016/j.eng.2019.12.012}. 

One of the primary challenges in developing effective threat models for visual adversarial attacks is the complexity and variability of the attacks themselves. As noted in \cite{doi:10.24963/ijcai.2021/694}, adversarial examples can be crafted to fool state-of-the-art image classifiers, even in the physical world. This highlights the need for threat models that can account for the diverse range of potential attacks and attackers. For instance, \cite{arxiv:1707.08945} demonstrates the feasibility of generating robust physical adversarial perturbations that can survive various physical transformations, making them a significant threat to real-world systems.

To address these challenges, researchers have proposed various threat models, including white-box, black-box, and gray-box models \cite{doi:10.1109/access.2018.2807385}. White-box models assume that the attacker has full knowledge of the system's architecture and parameters, while black-box models assume that the attacker has no knowledge of the system's internal workings. Gray-box models, on the other hand, assume that the attacker has some knowledge of the system's architecture and parameters, but not complete knowledge. Each of these models has its strengths and weaknesses, and the choice of model depends on the specific use case and the goals of the attacker \cite{doi:10.1007/s11633-019-1211-x}.

In addition to these models, researchers have also proposed various attack vectors, including patch-based, object-based, and scene-based attacks \cite{doi:10.1007/978-3-030-10925-7_4}. Patch-based attacks involve adding a small perturbation to the input image, while object-based attacks involve modifying the objects in the scene. Scene-based attacks, on the other hand, involve modifying the entire scene, such as by changing the lighting or adding new objects. Each of these attack vectors has its own strengths and weaknesses, and the choice of vector depends on the specific goals of the attacker and the vulnerabilities of the system \cite{arxiv:1802.06430}.

The development of effective threat models and attack vectors is crucial for defending against visual adversarial attacks. By understanding the potential attacks and vulnerabilities of a system, defenders can develop countermeasures that can mitigate the risks associated with these attacks. For instance, \cite{arxiv:2101.07922} proposes a novel approach to protecting social media users from facial recognition by leveraging adversarial attacks. Similarly, \cite{doi:10.1145/3433210.3437513} proposes a framework for constrained adversarial machine learning that can be used to defend against visual adversarial attacks in cyber-physical systems.

In conclusion, threat models play a vital role in understanding and defending against visual adversarial attacks. By providing a framework for analyzing the potential attacks and vulnerabilities of a system, threat models can help defenders develop effective countermeasures. The development of effective threat models and attack vectors is an active area of research, with various models and vectors being proposed to address the complexities and variability of visual adversarial attacks. As noted in \cite{doi:10.1145/3459637.3482029}, the study of adversarial robustness is crucial for developing trustworthy and reliable deep learning systems. Future research directions include the development of more sophisticated threat models that can account for the complexities of real-world systems and the development of more effective countermeasures that can mitigate the risks associated with visual adversarial attacks \cite{doi:10.1109/jiot.2021.3099164}.


\subsection{Types of Threat Models}


The scope of this subsection is to delve into the various types of threat models employed in visual adversarial attacks, with a particular emphasis on their applicability to physical world scenarios. As discussed in the previous section, threat models are essential in understanding and defending against visual adversarial attacks, as they outline the assumptions made about the attacker's knowledge, capabilities, and goals. The primary types of threat models include white-box, black-box, and gray-box attacks, each with its strengths, weaknesses, and suitability for different scenarios.

Building on the concept of threat models, it is crucial to examine the different types of threat models and their applicability to physical world scenarios. White-box threat models assume that the attacker has complete knowledge of the target system, including its architecture, parameters, and training data \cite{arxiv:1810.00069}. This level of access allows for highly targeted and effective attacks, as demonstrated by \cite{arxiv:1707.07397} and \cite{arxiv:1707.08945}. However, the assumption of full knowledge is often unrealistic in practical scenarios, limiting the applicability of white-box models.

In contrast, black-box threat models operate under the assumption that the attacker has no knowledge of the target system's internal workings \cite{doi:10.1049/cit2.12028}. This scenario is more realistic, as attackers often do not have access to the system's architecture or parameters. Black-box attacks, such as those presented in \cite{doi:10.1109/cvpr.2019.00790}, rely on querying the system to observe its behavior and craft adversarial examples. While black-box models are more realistic, they can be less effective than white-box models due to the lack of information about the system.

Gray-box threat models strike a balance between white-box and black-box models, assuming that the attacker has partial knowledge of the target system \cite{doi:10.1016/j.eng.2019.12.012}. This partial knowledge can include information about the system's architecture, training data, or parameters. Gray-box models are more flexible and can be used to simulate a variety of attack scenarios, making them useful for evaluating the robustness of systems \cite{doi:10.1109/access.2018.2807385}.

The choice of threat model depends on the specific use case and the level of access an attacker is assumed to have. In physical world scenarios, black-box and gray-box models are often more relevant, as attackers may not have direct access to the system's internal workings. However, white-box models can still be useful for evaluating the worst-case scenario and identifying potential vulnerabilities. As will be discussed in the following section, the selection of threat models has significant implications for the development of effective defense strategies against visual adversarial attacks.

Recent studies have also explored the use of threat models in specific domains, such as autonomous vehicles \cite{doi:10.1109/jiot.2021.3099164} and face recognition systems \cite{doi:10.1016/j.patcog.2022.109009}. These domain-specific threat models can provide valuable insights into the potential risks and vulnerabilities of visual adversarial attacks in real-world applications.

In conclusion, understanding the different types of threat models is crucial for developing effective defenses against visual adversarial attacks. By recognizing the strengths and weaknesses of each threat model, researchers and practitioners can design more robust systems that are better equipped to withstand attacks in various scenarios. As the field continues to evolve, it is essential to consider the applicability of threat models to physical world scenarios and to develop new models that can address the unique challenges of real-world attacks \cite{doi:10.1109/cvpr42600.2020.00108}. Future research should focus on developing more realistic and comprehensive threat models that can capture the complexities of physical world scenarios, ultimately leading to more secure and robust visual recognition systems.


\subsection{Attack Vectors}


The realm of visual adversarial attacks has witnessed a proliferation of innovative methods aimed at deceiving deep learning models, with various attack vectors being explored to compromise the security and reliability of these systems. Among these, patch-based, object-based, and scene-based attacks have garnered significant attention due to their potential to cause misclassification in different environments \cite{doi:10.1145/3128572.3140444}. Patch-based attacks, for instance, involve generating adversarial patches that can be applied to objects or scenes to deceive detectors \cite{arxiv:1707.08945}. These patches are often designed to be robust against various transformations and perturbations, ensuring their effectiveness in real-world scenarios \cite{doi:10.1007/978-3-030-10925-7_4}.

Object-based attacks, on the other hand, focus on modifying objects themselves to create adversarial examples \cite{arxiv:1807.07769}. This can be achieved through techniques such as 3D printing or adding stickers to objects, which can then be used to deceive object detectors \cite{doi:10.1109/iccvw.2019.00257}. The efficacy of these attacks has been demonstrated in various studies, including \cite{arxiv:1910.11099}, which showed that adversarial T-shirts can be designed to evade person detectors. Scene-based attacks, meanwhile, involve modifying the environment or scene in which objects are detected \cite{doi:10.1145/3377811.3380422}. This can be achieved through techniques such as projecting adversarial patterns onto surfaces or using adversarial billboards to deceive detectors.

The effectiveness of these attack vectors can be attributed to their ability to exploit the vulnerabilities of deep learning models, which are often designed to recognize patterns in data rather than understand the underlying context \cite{doi:10.1145/3274895.3274904}. As a result, adversarial examples can be crafted to mimic the patterns recognized by these models, leading to misclassification \cite{doi:10.1145/3398394}. The development of more sophisticated attack vectors, such as those that utilize 3D modeling and simulation \cite{doi:10.1016/j.neucom.2022.05.031}, has further exacerbated the problem, highlighting the need for more robust defense mechanisms.

In response to these threats, researchers have proposed various defense strategies, including adversarial training, input validation, and detection-based approaches \cite{arxiv:1712.07107}. Adversarial training involves training models on adversarial examples to improve their robustness, while input validation and detection-based approaches focus on identifying and mitigating the effects of adversarial attacks \cite{doi:10.1080/00031305.2021.2006781}. The development of more effective defense mechanisms, such as those that utilize generative adversarial networks \cite{doi:10.1109/access.2022.3160283}, has also been explored.

Despite these efforts, the cat-and-mouse game between attackers and defenders continues, with new attack vectors and defense mechanisms being developed in response to each other \cite{doi:10.1007/s11633-019-1211-x}. The use of techniques such as transfer learning and meta-learning has been proposed as a means of improving the robustness of models against adversarial attacks \cite{doi:10.1109/tip.2021.3092822}. Furthermore, the development of more sophisticated evaluation metrics and benchmarking protocols, such as those that incorporate human perception and environmental factors \cite{doi:10.1007/978-3-030-58589-1_3}, is necessary to accurately assess the effectiveness of defense mechanisms.

In conclusion, the various attack vectors used in visual adversarial attacks, including patch-based, object-based, and scene-based attacks, pose significant threats to the security and reliability of deep learning models. The development of more effective defense mechanisms, such as those that utilize adversarial training, input validation, and detection-based approaches, is necessary to mitigate these threats. Further research is needed to improve our understanding of the vulnerabilities of deep learning models and to develop more robust defense mechanisms against adversarial attacks \cite{doi:10.1109/biosig52210.2021.9548291}. By exploring new attack vectors and defense mechanisms, we can work towards developing more secure and reliable deep learning systems that can withstand the challenges posed by adversarial examples \cite{doi:10.1109/vr51125.2022.00073}. Ultimately, the development of more effective defense mechanisms will require a comprehensive understanding of the complex interplay between attack vectors, defense mechanisms, and the underlying vulnerabilities of deep learning models \cite{doi:10.1007/s13042-021-01435-0}.


\subsection{Physical Adversarial Attacks}


Physical adversarial attacks pose a significant threat to the security and reliability of computer vision systems in the physical world, as they involve manipulating real-world objects or environments to deceive these systems. These attacks have been shown to be effective against a wide range of applications, including object detection, facial recognition, and autonomous vehicles \cite{arxiv:1707.07397,arxiv:1707.08945}. A key challenge in generating physical adversarial examples is ensuring their robustness to various environmental conditions, such as lighting, viewpoint, and weather \cite{arxiv:1707.08945}. To address this challenge, researchers have proposed methods like the use of expectation over transformation (EOT) to improve the robustness of physical adversarial examples \cite{arxiv:1707.08945}.

The development of methods for generating physical adversarial examples that can be applied to real-world objects is also crucial. For instance, \cite{arxiv:1807.07769} proposed a method for generating physical adversarial examples to attack object detectors, while \cite{arxiv:1910.11099} demonstrated the use of adversarial T-shirts to evade person detectors. These methods have been effective in both digital and physical worlds \cite{arxiv:1807.07769,arxiv:1910.11099}. Furthermore, researchers have explored the use of various materials and techniques to create adversarial objects, such as customized stickers \cite{doi:10.1007/978-3-030-10925-7_4} and physical patches \cite{arxiv:1906.11897}, which have been shown to be effective in evading detection.

Evaluating physical adversarial attacks is another important area of research. Studies like \cite{doi:10.1007/978-3-030-58589-1_3} have proposed datasets for evaluating physical adversarial attacks on object detection, and \cite{doi:10.1109/iros51168.2021.9636638} has demonstrated the use of adversarial attacks to evaluate the robustness of camera-LiDAR models for 3D car detection. These evaluation methods are essential for understanding the vulnerabilities of computer vision systems to physical adversarial attacks \cite{doi:10.1007/978-3-030-58589-1_3,doi:10.1109/iros51168.2021.9636638}.

In response to the threats posed by physical adversarial attacks, researchers have proposed various defense strategies. For example, \cite{doi:10.1109/asp-dac47756.2020.9045584} proposed a comprehensive and lightweight defense methodology for detecting and mitigating physical adversarial attacks, and \cite{doi:10.1145/3450267.3450535} demonstrated the use of real-time detectors for the same purpose. These defense strategies have been shown to be effective in both digital and physical worlds \cite{doi:10.1109/asp-dac47756.2020.9045584,doi:10.1145/3450267.3450535}.

Overall, physical adversarial attacks pose a significant threat to computer vision systems, and addressing this threat requires the development of robust methods for generating and evaluating physical adversarial examples, as well as effective defense strategies. By understanding the challenges and opportunities in this area, researchers can work towards improving the security and reliability of computer vision systems in the physical world \cite{arxiv:1904.00759,arxiv:1910.11099,doi:10.1109/asp-dac47756.2020.9045584,doi:10.1145/3450267.3450535}.


\subsection{Defense Strategies}


The development of effective defense strategies against visual adversarial attacks is crucial for ensuring the reliability and security of computer vision systems in the physical world. Various approaches have been proposed to counter these attacks, including adversarial training, input preprocessing, and detection-based methods. Adversarial training involves training a model on a dataset that includes both normal and adversarial examples, with the goal of improving the model's robustness to attacks \cite{doi:10.1145/3128572.3140444}. This approach has been shown to be effective in improving the robustness of models against certain types of attacks, but its effectiveness can depend on the specific attack method used to generate the adversarial examples \cite{arxiv:1705.07204}.

Input preprocessing techniques, such as data normalization and feature extraction, can also be used to defend against adversarial attacks \cite{doi:10.1609/aaai.v32i1.11828}. These methods can help to reduce the sensitivity of models to small perturbations in the input data, making them more robust to attacks. However, the effectiveness of these methods can depend on the specific type of attack and the characteristics of the input data \cite{doi:10.1145/3274895.3274904}.

Detection-based approaches involve training a separate model to detect whether an input is an adversarial example or not \cite{doi:10.1109/tnnls.2021.3105238}. These methods can be effective in detecting certain types of attacks, but they can also be evaded by more sophisticated attack methods \cite{arxiv:1904.02884}. Recently, \cite{doi:10.1109/cvpr42600.2020.00108} proposed a novel approach to generate adversarial examples that are more natural and less detectable.

Another approach to defend against adversarial attacks is to use robust optimization techniques, such as minimax optimization, to improve the robustness of models \cite{arxiv:1901.08573}. These techniques can help to improve the robustness of models by optimizing the model's parameters to minimize the maximum loss over a set of possible inputs. However, the effectiveness of these techniques can depend on the specific type of attack and the characteristics of the input data \cite{arxiv:1904.13000}.

In addition to these approaches, several other methods have been proposed to defend against adversarial attacks, including the use of ensemble methods, which involve combining the predictions of multiple models to improve robustness \cite{arxiv:1705.07204}. These methods can be effective in improving the robustness of models, but they can also increase the computational cost of the model \cite{doi:10.1109/asp-dac47756.2020.9045584}.

Recently, \cite{doi:10.1109/cvpr42600.2020.01426} proposed a novel approach to generate physical-world-resilient adversarial examples for autonomous driving. \cite{arxiv:2012.12235} proposed a framework to design objects that are robust to adversarial attacks. \cite{doi:10.1609/aaai.v35i9.16927} proposed a novel approach to improve the robustness of models to natural perturbations.

In conclusion, defending against visual adversarial attacks in the physical world is a complex task that requires a comprehensive approach. Various defense strategies have been proposed, including adversarial training, input preprocessing, detection-based methods, and robust optimization techniques. While these approaches have shown promise in improving the robustness of models, they can also have limitations and trade-offs. Further research is needed to develop more effective and efficient defense strategies that can be used in a variety of applications \cite{doi:10.24963/ijcai.2021/694}. As \cite{arxiv:2209.14262} pointed out, it is essential to develop a framework that can enable a comprehensive evaluation of the vulnerability of face recognition in the physical world. By combining different approaches and techniques, it may be possible to develop more robust and reliable computer vision systems that can withstand adversarial attacks in the physical world \cite{doi:10.1016/j.patcog.2022.109009}.


\subsection{Evaluation Metrics and Benchmarking}


The evaluation of adversarial robustness in deep learning models is a critical aspect of ensuring their reliability and security in the physical world, as it directly impacts the effectiveness of defense strategies against visual adversarial attacks. As discussed in the previous section, defending against visual adversarial attacks in the physical world is a complex task that requires a comprehensive approach, including the development of robust evaluation metrics and benchmarking protocols. Establishing robust evaluation metrics and benchmarking protocols is essential for assessing the effectiveness of these models against various types of attacks. As \cite{arxiv:1911.05268} notes, the development of comprehensive evaluation metrics is crucial for understanding the vulnerabilities of deep learning models. Existing metrics, such as accuracy, robustness, and attack success rate, have their strengths and limitations. For instance, \cite{doi:10.1016/j.eng.2019.12.012} highlights the importance of considering both the digital and physical aspects of adversarial attacks when evaluating the robustness of deep learning models.

One of the key challenges in evaluating adversarial robustness is the lack of standardized benchmarking protocols. As \cite{doi:10.1109/cvpr42600.2020.00108} suggests, benchmarking protocols should be designed to capture the variability of physical-world scenarios, including different lighting conditions, viewpoints, and environmental factors. The development of such protocols requires a comprehensive understanding of the physical world and its potential impact on deep learning models. \cite{doi:10.1109/cvpr42600.2020.01426} proposes a framework for generating physical-world-resilient adversarial examples, which can be used to evaluate the robustness of deep learning models in autonomous driving scenarios. This framework can be used in conjunction with defense strategies, such as adversarial training and input preprocessing, to improve the robustness of models against physical-world attacks.

The evaluation of adversarial robustness also requires consideration of the attack vectors and threat models used in physical-world attacks. As \cite{arxiv:1707.08945} notes, physical-world attacks can be more challenging to defend against than digital attacks, due to the complexity of the physical environment. \cite{arxiv:1712.09665} proposes a method for generating universal, robust, targeted adversarial image patches in the real world, which can be used to evaluate the robustness of deep learning models against physical-world attacks. Furthermore, the consideration of countermeasures against adversarial attacks, such as defense strategies, is crucial for developing effective evaluation metrics and benchmarking protocols.

Emerging trends in evaluation metrics and benchmarking include the use of more realistic and comprehensive metrics, such as those that capture the complexity of physical-world scenarios and the variability of adversarial attacks. As \cite{doi:10.1201/9781351251389-8} suggests, the development of such metrics requires a deep understanding of the physical world and its potential impact on deep learning models. \cite{doi:10.1007/978-3-030-58589-1_3} proposes a dataset of physical adversarial attacks on object detection, which can be used to evaluate the robustness of deep learning models against physical-world attacks. These emerging trends and datasets can be used to inform the development of more effective defense strategies and evaluation metrics, ultimately leading to more robust and reliable computer vision systems.

In conclusion, the evaluation of adversarial robustness in deep learning models is a critical aspect of ensuring their reliability and security in the physical world. The development of robust evaluation metrics and benchmarking protocols is essential for assessing the effectiveness of these models against various types of attacks. As \cite{arxiv:2104.02361} notes, the evaluation of adversarial robustness requires consideration of the physical world and its potential impact on deep learning models. Future research directions include the development of more realistic and comprehensive evaluation metrics, the consideration of emerging trends and challenges, and the exploration of innovative perspectives and solutions for defending against adversarial attacks, as discussed in \cite{doi:10.1109/iros51168.2021.9636638}, \cite{arxiv:2209.09502}, and \cite{doi:10.1109/eurosp53844.2022.00047}.


\section{Defenses Against Visual Adversarial Attacks}



\subsection{Introduction to Defense Mechanisms}


The rise of visual adversarial attacks has posed significant threats to the security and reliability of computer vision systems, emphasizing the need for robust defense mechanisms. As discussed in \cite{arxiv:1810.00069}, the development of effective defenses is crucial to mitigate the risks associated with these attacks. Defense mechanisms against visual adversarial attacks can be broadly categorized into several types, including adversarial training, input validation, and detection techniques. Adversarial training, as explored in \cite{arxiv:2102.01356}, involves training a model on adversarial examples to enhance its robustness. This approach has been shown to be effective in improving the model's resilience to attacks, but it can be computationally expensive and may not provide complete protection against all types of attacks.

Input validation techniques, on the other hand, aim to detect and prevent adversarial examples from being processed by the model. As discussed in \cite{doi:10.1109/asp-dac47756.2020.9045584}, these techniques can be used to filter out adversarial examples or to transform the input data to reduce the impact of the attack. However, the effectiveness of these techniques can be limited by the complexity of the attack and the quality of the input data. Detection techniques, as presented in \cite{arxiv:1906.11897}, involve identifying adversarial examples and preventing them from being processed by the model. These techniques can be based on statistical analysis, machine learning models, or other methods, and can be used in conjunction with other defense mechanisms to provide an additional layer of protection.

The choice of defense mechanism depends on the specific application, the type of attack, and the level of security required. As noted in \cite{doi:10.1109/iros51168.2021.9636638}, a combination of different defense mechanisms can provide more effective protection against visual adversarial attacks. Furthermore, the use of ensemble methods, as discussed in \cite{arxiv:1810.00069}, can also improve the robustness of computer vision systems against these attacks. In addition to these approaches, other defense strategies, such as \cite{doi:10.1145/3433210.3437513}, have been proposed to protect against physical adversarial attacks.

Recent studies, such as \cite{doi:10.1109/cvpr42600.2020.01426}, have also explored the use of generative models to generate physical-world-resilient adversarial examples. These examples can be used to test the robustness of computer vision systems and to develop more effective defense mechanisms. Moreover, the use of adversarial training and robust optimization techniques, as discussed in \cite{arxiv:1707.08945}, can also improve the robustness of computer vision systems against physical adversarial attacks.

In conclusion, the development of effective defense mechanisms against visual adversarial attacks is a critical area of research, with significant implications for the security and reliability of computer vision systems. As noted in \cite{doi:10.1145/3459637.3482029}, the design of robust defense mechanisms requires a deep understanding of the attack strategies and the vulnerabilities of computer vision systems. By combining different defense mechanisms and using ensemble methods, it is possible to provide more effective protection against visual adversarial attacks. Future research should focus on developing more effective and efficient defense mechanisms, as well as exploring new approaches, such as \cite{doi:10.1145/3433210.3437513}, to protect against these attacks. As discussed in \cite{doi:10.1109/jiot.2021.3099164}, the evaluation of the robustness of computer vision systems against visual adversarial attacks is also crucial to ensure the safety and reliability of these systems.


\subsection{Adversarial Training and Robust Optimization}


Adversarial training and robust optimization are two closely related techniques used to defend against visual adversarial attacks, which are a critical concern in the development of secure and reliable computer vision systems. As discussed in the previous section, the rise of visual adversarial attacks has posed significant threats to the security and reliability of computer vision systems, emphasizing the need for robust defense mechanisms. Adversarial training involves training a model on a dataset that includes adversarial examples, with the goal of improving the model's robustness to such attacks \cite{arxiv:1810.00069}. Robust optimization, on the other hand, refers to the use of optimization techniques that are designed to minimize the worst-case loss of a model, rather than the average loss \cite{arxiv:1707.08945}. One popular robust optimization technique is minimax optimization, which involves minimizing the maximum loss of a model over a set of possible inputs \cite{doi:10.1201/9781351251389-8}.

The use of adversarial training and robust optimization techniques has been shown to be effective in defending against visual adversarial attacks. For example, \cite{arxiv:1707.07397} demonstrates the use of adversarial training to defend against physical adversarial attacks on object detectors. Similarly, \cite{arxiv:1707.08945} shows the effectiveness of robust optimization techniques in defending against physical adversarial attacks on deep learning models. Other works, such as \cite{doi:10.1126/science.aaw4399} and \cite{arxiv:1810.00069}, also highlight the importance of using adversarial training and robust optimization techniques to defend against visual adversarial attacks. These techniques can be used in conjunction with other defense mechanisms, such as input validation and detection techniques, to provide a multi-layered defense strategy against visual adversarial attacks.

One of the key challenges in using adversarial training and robust optimization techniques is the need to balance the trade-off between model robustness and accuracy. As noted in \cite{doi:10.1049/cit2.12028}, increasing the robustness of a model can often come at the cost of decreasing its accuracy. Therefore, it is essential to carefully tune the parameters of the adversarial training and robust optimization algorithms to achieve the desired balance between robustness and accuracy. Additionally, \cite{doi:10.1109/cvpr42600.2020.00108} highlights the importance of considering the physical-world constraints when designing adversarial training and robust optimization techniques. By taking into account the physical-world constraints and carefully tuning the parameters of these techniques, it is possible to develop effective defense methods that can protect against a wide range of visual adversarial attacks.

Recent advances in adversarial training and robust optimization techniques have also led to the development of more sophisticated defense methods. For example, \cite{doi:10.1109/cvpr42600.2020.01426} proposes the use of generative adversarial networks (GANs) to generate physical-world-resilient adversarial examples for autonomous driving systems. Similarly, \cite{doi:10.1109/asp-dac47756.2020.9045584} proposes a comprehensive and lightweight defense methodology against physical adversarial attacks on embedded multimedia applications. These advances have significant implications for the development of secure and reliable computer vision systems, and highlight the importance of continued research in this area.

In conclusion, adversarial training and robust optimization techniques are essential components of any defense strategy against visual adversarial attacks. By carefully tuning the parameters of these techniques and considering the physical-world constraints, it is possible to develop effective defense methods that can protect against a wide range of visual adversarial attacks. As noted in \cite{doi:10.1109/access.2022.3208131}, the development of more sophisticated defense methods will require continued advances in adversarial training and robust optimization techniques. The use of these techniques, in conjunction with other defense mechanisms such as input validation and preprocessing techniques, will be crucial in developing secure and reliable computer vision systems. Future research directions may include the development of more efficient and effective adversarial training and robust optimization algorithms, as well as the exploration of new defense methods that can be used in conjunction with these techniques \cite{doi:10.1155/2021/4907754}.


\subsection{Input Validation and Preprocessing}


Input validation and preprocessing techniques have emerged as crucial defenses against visual adversarial attacks, which can compromise the security and reliability of computer vision systems. These techniques aim to detect and prevent adversarial attacks by sanitizing and normalizing the input data, thereby reducing the vulnerability of deep learning models to such attacks. As highlighted in \cite{doi:10.1145/3128572.3140444}, the development of effective input validation and preprocessing techniques is essential to counter the growing threat of adversarial examples.

One of the primary approaches to input validation is input sanitization, which involves removing or modifying potentially malicious components of the input data. This can be achieved through techniques such as data normalization, feature scaling, and outlier detection. For instance, \cite{arxiv:1707.08945} demonstrates the effectiveness of input sanitization in defending against physical-world adversarial attacks. The authors propose a method for generating robust physical adversarial examples that can withstand various environmental conditions, and show that input sanitization can significantly improve the robustness of deep learning models against such attacks.

Another important technique for input validation is input normalization, which involves transforming the input data into a standard format to reduce the impact of adversarial perturbations. This can be achieved through techniques such as min-max scaling, standardization, and normalization. As shown in \cite{doi:10.1109/tip.2021.3092822}, input normalization can be an effective defense against black-box adversarial attacks, which are particularly challenging to detect and prevent.

In addition to input sanitization and normalization, other techniques such as data augmentation and feature extraction can also be used to improve the robustness of deep learning models against adversarial attacks. For example, \cite{doi:10.1109/iccvw.2019.00257} proposes a method for generating adversarial examples using a generative adversarial network (GAN) framework, and demonstrates the effectiveness of data augmentation in improving the robustness of deep learning models against such attacks.

The use of input validation and preprocessing techniques can also be combined with other defense mechanisms, such as adversarial training and detection-based approaches, to provide a more comprehensive defense against visual adversarial attacks. As highlighted in \cite{arxiv:1712.07107}, the combination of input validation and adversarial training can provide a robust defense against adversarial attacks, and can significantly improve the security and reliability of computer vision systems.

However, despite the effectiveness of input validation and preprocessing techniques, there are also several challenges and limitations associated with these approaches. For instance, \cite{arxiv:1711.08478} demonstrates that some input validation techniques can be bypassed by sophisticated adversarial attacks, and highlights the need for more robust and effective defense mechanisms.

In conclusion, input validation and preprocessing techniques are essential defenses against visual adversarial attacks, and can significantly improve the security and reliability of computer vision systems. The development of effective input validation and preprocessing techniques requires a deep understanding of the underlying mechanisms of adversarial attacks, as well as the strengths and limitations of different defense mechanisms. As highlighted in \cite{doi:10.1186/s42400-018-0012-9}, the study of input validation and preprocessing techniques is an active area of research, and there are many opportunities for future work and innovation in this field. Further research is needed to develop more robust and effective defense mechanisms, and to address the challenges and limitations associated with input validation and preprocessing techniques.


\subsection{Detection-Based Approaches}


Detection-based approaches have emerged as a crucial defense mechanism against visual adversarial attacks, aiming to identify and prevent such attacks from compromising the integrity of computer vision systems. As a complementary approach to input validation and preprocessing techniques, detection-based methods leverage various detection techniques, including anomaly detection and classification, to distinguish between legitimate and adversarial inputs \cite{doi:10.1201/9781351251389-8}. By analyzing the characteristics of adversarial examples, detection-based methods can effectively flag potentially harmful inputs, thereby safeguarding the system against malicious attacks \cite{arxiv:1707.08945}.

One of the primary advantages of detection-based approaches is their ability to operate independently of the specific attack type or model architecture \cite{doi:10.1145/3319535.3339815}. This flexibility allows them to be applied in a wide range of scenarios, from image classification to object detection, and even in complex settings such as autonomous driving \cite{doi:10.1609/aaai.v35i2.16211}. For instance, \cite{arxiv:1906.11897} demonstrates the effectiveness of detection-based methods in identifying adversarial patches designed to deceive object detectors. Furthermore, detection-based approaches can be used in conjunction with input validation and preprocessing techniques to provide a more robust defense against visual adversarial attacks.

Anomaly detection is a key component of detection-based approaches, as it enables the identification of inputs that deviate from the expected norm \cite{doi:10.1145/3450267.3450535}. By training on a dataset of legitimate inputs, anomaly detection models can learn to recognize patterns that are indicative of adversarial examples \cite{arxiv:1910.11099}. Classification-based detection methods, on the other hand, rely on labeled datasets to learn the differences between legitimate and adversarial inputs \cite{doi:10.1007/978-3-030-58548-8_1}. The combination of anomaly detection and classification-based methods can provide a more comprehensive defense against visual adversarial attacks.

The strengths and limitations of detection-based approaches have been extensively evaluated in various studies \cite{doi:10.1109/iros51168.2021.9636638,doi:10.1109/asp-dac47756.2020.9045584}. While these methods have shown promise in detecting adversarial examples, they can be computationally expensive and may require significant amounts of labeled data \cite{doi:10.1609/aaai.v35i12.17259}. Furthermore, the effectiveness of detection-based approaches can be compromised by the sophistication of the attack, with more complex attacks potentially evading detection \cite{doi:10.1109/cvpr46437.2021.00614}. Therefore, it is essential to continue advancing the field of detection-based approaches to address these challenges and limitations.

Emerging trends in detection-based approaches include the use of ensemble methods, which combine multiple detection models to improve overall performance \cite{arxiv:2110.03825}. Additionally, the integration of physical-world constraints, such as lighting and viewpoint, can enhance the robustness of detection-based methods \cite{arxiv:2007.16118}. These advancements can be used to inform the development of more efficient and effective detection methods, as well as explore the application of detection-based approaches in diverse domains, such as autonomous driving and surveillance systems \cite{doi:10.1007/978-3-030-58589-1_3}. The future of detection-based approaches holds much promise, and their continued development will be crucial in defending against visual adversarial attacks.

In conclusion, detection-based approaches offer a promising defense mechanism against visual adversarial attacks, with various techniques and methods being developed to improve their effectiveness. By understanding the strengths and limitations of these approaches, researchers can continue to advance the field, ultimately leading to more robust and secure computer vision systems \cite{doi:10.1109/wacv51458.2022.00144,doi:10.1109/ijcnn55064.2022.9892612}. As the field continues to evolve, it is essential to consider the potential applications and implications of detection-based approaches, as well as the challenges that must be addressed to ensure their successful deployment in real-world scenarios \cite{doi:10.1109/eurosp53844.2022.00047}. The development of ensemble methods and hybrid approaches, which will be discussed in the following section, can provide a more comprehensive defense against visual adversarial attacks, and their combination with detection-based approaches can lead to even more robust and secure computer vision systems.


\subsection{Ensemble Methods and Hybrid Approaches}


Ensemble methods and hybrid approaches have emerged as a promising direction in defending against visual adversarial attacks. The core idea behind these methods is to combine the strengths of multiple models or techniques to improve the overall robustness against adversarial examples. One of the most popular ensemble methods is bagging, which involves training multiple models on different subsets of the training data and then combining their predictions \cite{doi:10.1145/3128572.3140444}. Another approach is boosting, which iteratively trains models on the mistakes made by previous models, resulting in a strong ensemble that is more robust to adversarial attacks \cite{arxiv:1705.07204}.

Hybrid approaches, on the other hand, combine different techniques, such as adversarial training and input validation, to improve robustness. For example, \cite{arxiv:1901.08573} proposes a hybrid approach that combines adversarial training with a regularization term to encourage the model to produce more robust features. Similarly, \cite{arxiv:1902.06705} suggests a hybrid approach that combines multiple models with different architectures and training procedures to improve overall robustness.

One of the key benefits of ensemble methods and hybrid approaches is that they can provide a more comprehensive defense against visual adversarial attacks. By combining multiple models or techniques, these approaches can reduce the risk of overfitting to specific types of attacks and improve the overall robustness of the system \cite{doi:10.1201/9781351251389-8}. Additionally, ensemble methods and hybrid approaches can be used to improve the transferability of adversarial examples, making it more difficult for attackers to craft effective attacks \cite{arxiv:1904.02884}.

However, ensemble methods and hybrid approaches also have some limitations and challenges. One of the main challenges is the increased computational cost and complexity of these approaches, which can make them more difficult to implement and train \cite{doi:10.1007/978-3-030-01258-8_39}. Additionally, the choice of models or techniques to combine can have a significant impact on the overall performance of the ensemble, and finding the optimal combination can be a challenging task \cite{arxiv:1904.13000}.

Recent studies have also explored the use of ensemble methods and hybrid approaches in the context of physical-world attacks. For example, \cite{doi:10.1007/978-3-030-58589-1_3} proposes a hybrid approach that combines adversarial training with a physical-world attack simulator to improve the robustness of object detection models against physical-world attacks. Similarly, \cite{arxiv:1707.08945} suggests a hybrid approach that combines multiple models with different architectures and training procedures to improve the overall robustness of deep learning models against physical-world attacks.

In conclusion, ensemble methods and hybrid approaches have shown great promise in defending against visual adversarial attacks. By combining the strengths of multiple models or techniques, these approaches can provide a more comprehensive defense against adversarial examples and improve the overall robustness of the system. However, these approaches also have some limitations and challenges, and further research is needed to fully explore their potential and develop more effective and efficient defenses against visual adversarial attacks \cite{arxiv:1911.05268}. Future directions for research include exploring new ensemble methods and hybrid approaches, improving the transferability of adversarial examples, and developing more effective and efficient defenses against physical-world attacks \cite{arxiv:2209.14262}.


\subsection{Evaluation Metrics and Benchmarking}


Evaluating the effectiveness of defense mechanisms against visual adversarial attacks is crucial for ensuring the security and reliability of deep learning models in the physical world. As discussed in the previous section, ensemble methods and hybrid approaches have shown great promise in defending against visual adversarial attacks. However, to assess the robustness of these models, various evaluation metrics and benchmarking protocols have been proposed \cite{arxiv:1712.07107}. One of the primary metrics used to evaluate the performance of defense mechanisms is accuracy, which measures the proportion of correctly classified samples \cite{doi:10.1016/j.eng.2019.12.012}. However, accuracy alone is not sufficient to evaluate the robustness of a model, as it does not account for the model's ability to withstand adversarial attacks.

To provide a more comprehensive understanding of a model's robustness, robustness metrics such as the adversarial loss and the attack success rate are more suitable for evaluating the effectiveness of defense mechanisms \cite{doi:10.1145/3274895.3274904}. The adversarial loss measures the difference between the model's predictions on clean and adversarial samples, while the attack success rate measures the proportion of samples that are misclassified by the model when subjected to adversarial attacks \cite{doi:10.1109/iros51168.2021.9636638}. These metrics can be used to compare the effectiveness of different defense mechanisms, including ensemble methods and hybrid approaches. Benchmarking protocols are also essential for evaluating the performance of defense mechanisms \cite{doi:10.1109/access.2018.2807385}. These protocols provide a standardized framework for evaluating the robustness of models against various types of adversarial attacks \cite{arxiv:1712.09665}.

Examples of benchmarking protocols include the APRICOT dataset, which provides a collection of physical adversarial patches that can be used to evaluate the robustness of object detection models \cite{doi:10.1007/978-3-030-58589-1_3}, and the CARLA simulator, which provides a platform for evaluating the robustness of autonomous driving models against physical adversarial attacks \cite{arxiv:2007.16118}. In addition to these metrics and benchmarking protocols, other evaluation methods have been proposed, such as the use of adversarial training and robust optimization techniques \cite{arxiv:2102.01356}. Adversarial training involves training a model on a mixture of clean and adversarial samples, while robust optimization techniques involve optimizing the model's parameters to minimize the adversarial loss \cite{doi:10.1609/aaai.v34i01.5443}.

Despite these advances, evaluating the effectiveness of defense mechanisms against visual adversarial attacks remains a challenging task \cite{arxiv:1810.00069}. The complexity of the physical world, the variability of lighting and environmental conditions, and the diversity of potential attack scenarios all contribute to the difficulty of evaluating the robustness of models \cite{doi:10.1109/lcomm.2019.2901469}. Furthermore, the lack of standardized evaluation metrics and benchmarking protocols makes it difficult to compare the effectiveness of different defense mechanisms \cite{doi:10.1007/s11633-019-1211-x}. To address these challenges, future research should focus on developing more comprehensive and realistic evaluation metrics and benchmarking protocols \cite{doi:10.1155/2021/4907754}. These protocols should account for the variability of the physical world and the diversity of potential attack scenarios \cite{doi:10.1145/3433210.3437513}.

The development of new evaluation metrics, such as those that incorporate human perception and environmental factors, is also crucial for advancing the field \cite{doi:10.1007/s13042-021-01435-0}. These metrics can provide a more comprehensive understanding of the robustness of models and can be used to develop more effective defense mechanisms \cite{doi:10.1109/wacvw54805.2022.00036}. Furthermore, the use of generative models, such as GANs, can be explored for generating physical adversarial examples that are more realistic and effective \cite{arxiv:2209.09502}. By advancing the field of evaluation metrics and benchmarking protocols, we can improve the security and reliability of deep learning models in the physical world and ensure their safe and effective deployment in a wide range of applications \cite{doi:10.32604/jai.2021.014175}. Ultimately, the development of comprehensive and realistic evaluation metrics and benchmarking protocols will play a critical role in the development of effective defense mechanisms against visual adversarial attacks, which will be discussed in the next section.


\section{Applications and Case Studies of Visual Adversarial Attacks}



\subsection{Introduction to Applications of Visual Adversarial Attacks}


The application of visual adversarial attacks in real-world scenarios has significant implications for various domains, including security, safety, and reliability of computer vision systems. As highlighted in \cite{doi:10.24963/ijcai.2021/694}, the vulnerability of deep neural networks to adversarial examples poses a substantial threat to the deployment of these systems in safety-critical applications. The potential impact of visual adversarial attacks on autonomous vehicles, face recognition systems, and other computer vision-based applications cannot be overstated. For instance, \cite{arxiv:1707.08945} demonstrates the feasibility of generating robust physical adversarial perturbations that can mislead object detectors in the physical world.

Recent studies have explored the application of visual adversarial attacks in various domains, including autonomous driving \cite{doi:10.1016/j.sysarc.2020.101766}, face recognition \cite{doi:10.1609/aaai.v32i1.12341}, and aerial detection \cite{doi:10.1109/tgrs.2022.3225306}. These studies have shown that visual adversarial attacks can be used to compromise the security and reliability of computer vision systems, highlighting the need for robust defense mechanisms. As noted in \cite{doi:10.1016/j.eng.2019.12.012}, the development of effective defense strategies against visual adversarial attacks is crucial to ensuring the reliability and security of computer vision systems.

The concept of physical adversarial attacks, which involves manipulating real-world objects or environments to deceive computer vision systems, has gained significant attention in recent years. \cite{doi:10.1007/978-3-030-10925-7_4} proposes a method for generating robust physical adversarial examples that can evade object detection systems. Similarly, \cite{doi:10.1109/cvpr42600.2020.01426} presents a framework for generating physical-world-resilient adversarial examples for autonomous driving scenarios. These studies demonstrate the potential of physical adversarial attacks to compromise the security and reliability of computer vision systems in the physical world.

The evaluation of visual adversarial attacks in real-world scenarios is a challenging task, requiring the development of robust evaluation metrics and benchmarking protocols. As highlighted in \cite{doi:10.1109/jiot.2021.3099164}, the assessment of driving safety in vision-based autonomous vehicles is crucial to ensuring the reliability and security of these systems. The development of standardized evaluation protocols and benchmarking datasets for visual adversarial attacks is essential to facilitating comparison and advancement in the field. \cite{doi:10.1007/978-3-030-58589-1_3} presents a dataset of physical adversarial attacks on object detection, providing a benchmark for evaluating the effectiveness of defense mechanisms against visual adversarial attacks.

In conclusion, the application of visual adversarial attacks in real-world scenarios has significant implications for various domains, including security, safety, and reliability of computer vision systems. The development of robust defense mechanisms against visual adversarial attacks is crucial to ensuring the reliability and security of these systems. As noted in \cite{doi:10.1145/3459637.3482029}, the development of effective defense strategies against visual adversarial attacks requires a comprehensive understanding of the underlying principles and mechanisms. Future research directions include the development of standardized evaluation protocols and benchmarking datasets, as well as the exploration of novel defense mechanisms against visual adversarial attacks. As highlighted in \cite{arxiv:2101.07922}, the potential applications of visual adversarial attacks extend beyond security and safety to include privacy protection and data security. Overall, the study of visual adversarial attacks has significant implications for the development of reliable and secure computer vision systems, and ongoing research in this area is essential to addressing the challenges and opportunities presented by these attacks.


\subsection{Case Studies of Visual Adversarial Attacks on Autonomous Vehicles}


The security and reliability of autonomous vehicles are of paramount importance, given their potential to revolutionize transportation systems. As discussed in the previous section, the application of visual adversarial attacks in real-world scenarios has significant implications for various domains, including security, safety, and reliability of computer vision systems. Recent studies have shown that autonomous vehicles are vulnerable to visual adversarial attacks, which can compromise their safety and reliability \cite{arxiv:1707.08945}. These attacks involve crafting adversarial examples that can deceive the vehicle's perception system, leading to misclassification or misdetection of objects \cite{doi:10.24963/ijcai.2021/694}.

Several studies have demonstrated the feasibility of generating robust physical adversarial examples that can attack autonomous vehicles. For instance, \cite{arxiv:1707.07397} demonstrated the existence of robust 3D adversarial objects that can fool image classifiers under various real-world conditions. The study showed that these objects can be synthesized using a combination of printing and projecting techniques, and can be used to create physical adversarial examples that can deceive autonomous vehicles. Similarly, \cite{arxiv:1707.08945} proposed a method for generating robust physical adversarial examples that can attack object detectors in the physical world. The study demonstrated the effectiveness of the proposed method using a stop sign as an example, and showed that the generated adversarial examples can be used to fool state-of-the-art object detectors.

Other studies have also explored the vulnerability of autonomous vehicles to visual adversarial attacks. \cite{doi:10.1007/978-3-030-10925-7_4} proposed a method for generating physical adversarial examples that can attack Faster R-CNN object detectors. The study demonstrated the effectiveness of the proposed method using a stop sign as an example, and showed that the generated adversarial examples can be used to fool state-of-the-art object detectors. Additionally, \cite{doi:10.1109/cvpr42600.2020.01426} proposed a method for generating physical-world-resilient adversarial examples for autonomous driving. The study demonstrated the effectiveness of the proposed method using a variety of examples, and showed that the generated adversarial examples can be used to fool state-of-the-art object detectors.

The potential risks and consequences of visual adversarial attacks on autonomous vehicles are significant. The generated adversarial examples can be used to deceive the vehicle's perception system, leading to misclassification or misdetection of objects. This can have serious consequences, including accidents and injuries. Therefore, it is essential to develop effective defense mechanisms against these attacks. \cite{doi:10.1109/asp-dac47756.2020.9045584} proposed a comprehensive and lightweight CNN defense methodology against physical adversarial attacks on embedded multimedia applications. The study demonstrated the effectiveness of the proposed method using a variety of examples, and showed that the proposed method can be used to defend against physical adversarial attacks.

In conclusion, visual adversarial attacks on autonomous vehicles pose a significant threat to their safety and reliability. The generated adversarial examples can be used to deceive the vehicle's perception system, leading to misclassification or misdetection of objects. As we move forward to the next section, which discusses the application of visual adversarial attacks on face recognition systems, it is essential to note that the development of effective defense mechanisms against these attacks is crucial to ensuring the reliability and security of computer vision systems. Future research should focus on developing more effective defense mechanisms, as well as investigating the potential risks and consequences of these attacks. As noted by \cite{doi:10.1109/iros51168.2021.9636638}, the vulnerability of autonomous vehicles to visual adversarial attacks is a critical issue that needs to be addressed. By developing effective defense mechanisms and investigating the potential risks and consequences of these attacks, we can ensure the safe and reliable operation of autonomous vehicles.


\subsection{Visual Adversarial Attacks on Face Recognition Systems}


The application of visual adversarial attacks on face recognition systems has garnered significant attention in recent years, owing to the potential risks and consequences of compromising the security and reliability of these systems. Face recognition technology is widely used in various domains, including security, law enforcement, and border control, making it an attractive target for adversarial attacks. Researchers have demonstrated the feasibility of generating adversarial examples that can fool face recognition systems, both in digital and physical domains. For instance, \cite{arxiv:1910.11099} have shown that adversarial T-shirts can be designed to evade face detection and recognition systems.

The generation of adversarial examples for face recognition systems typically involves optimizing a perturbation to be added to the input image, such that the resulting image is misclassified by the system. Various methods have been proposed to generate such adversarial examples, including \cite{doi:10.1109/iccvw.2019.00257} and \cite{doi:10.32604/jai.2021.014175}. These methods often rely on gradient-based optimization techniques to find the optimal perturbation. However, the effectiveness of these methods can be limited by the robustness of the face recognition system and the quality of the input images.

To improve the robustness of face recognition systems against adversarial attacks, various defense mechanisms have been proposed. These include \cite{doi:10.1109/asp-dac47756.2020.9045584} and \cite{doi:10.1109/eurosp53844.2022.00047}. Adversarial training involves training the face recognition system on a dataset that includes adversarial examples, with the goal of improving its robustness to such attacks. Input validation and preprocessing techniques aim to detect and remove adversarial perturbations from the input images. Detection-based approaches, on the other hand, focus on identifying adversarial examples at test time, using techniques such as anomaly detection and classification.

Despite these efforts, the development of effective defense mechanisms against visual adversarial attacks on face recognition systems remains an open challenge. \cite{doi:10.1186/s42400-018-0012-9} highlights the need for more research on the generation and detection of adversarial examples in the physical domain. \cite{arxiv:1707.08945} demonstrates the feasibility of generating physical adversarial examples that can fool face recognition systems in the wild. These examples can be created using various methods, including \cite{doi:10.1016/j.patcog.2022.109009} and \cite{doi:10.1007/978-3-030-58589-1_3}.

The implications of visual adversarial attacks on face recognition systems are far-reaching and concerning. \cite{doi:10.1201/9781351251389-8} highlights the potential risks of such attacks, including the compromise of national security and the erosion of trust in face recognition technology. \cite{doi:10.1109/tip.2021.3092822} demonstrates the feasibility of generating adversarial examples that can fool face recognition systems in a black-box setting, without access to the system's internal workings. \cite{doi:10.1016/j.patcog.2022.109009} shows that physical adversarial examples can be generated to fool face recognition systems in the wild, using a combination of digital and physical attacks.

In conclusion, visual adversarial attacks on face recognition systems pose a significant threat to the security and reliability of these systems. The generation of adversarial examples, both in digital and physical domains, has been demonstrated to be feasible using various methods. The development of effective defense mechanisms against such attacks remains an open challenge, requiring further research and innovation. As face recognition technology continues to be widely adopted, it is essential to address the vulnerabilities of these systems to adversarial attacks, to ensure their safe and reliable operation. Future research directions include the development of more robust face recognition systems, the creation of more effective defense mechanisms, and the investigation of the potential consequences of visual adversarial attacks on face recognition systems. \cite{arxiv:1904.00759} and \cite{doi:10.1109/vr51125.2022.00073} provide interesting insights into the potential of projector-based attacks and camera-based attacks, respectively.


\subsection{Defenses Against Visual Adversarial Attacks in the Physical World}


Defenses against visual adversarial attacks in the physical world are crucial for ensuring the security and reliability of computer vision systems. As \cite{arxiv:1707.08945} and \cite{doi:10.1109/lcomm.2019.2901469} have shown, physical adversarial attacks can be highly effective in compromising the accuracy of deep learning models. To counter these threats, various defense strategies have been developed, including adversarial training, input validation, and detection-based approaches. These defense strategies are essential for mitigating the risks associated with visual adversarial attacks, which can have significant consequences, as highlighted in the previous section on the vulnerability of face recognition systems to adversarial attacks.

Adversarial training involves training a model on a dataset that includes both clean and adversarial examples, with the goal of improving the model's robustness to attacks. This approach has been shown to be effective in defending against digital adversarial attacks, but its applicability to physical attacks is still limited. For instance, \cite{doi:10.1109/cvpr42600.2020.01426} proposed a method for generating physical-world-resilient adversarial examples, but the effectiveness of this approach in defending against physical attacks is still uncertain. Furthermore, the development of more robust face recognition systems, as discussed in the previous section, can also benefit from adversarial training.

Input validation techniques, such as input sanitization and normalization, can also be used to defend against physical adversarial attacks. These methods aim to detect and remove adversarial perturbations from the input data, but they may not be effective against sophisticated attacks that can evade detection. For example, \cite{doi:10.1109/asp-dac47756.2020.9045584} proposed a lightweight defense methodology that uses input validation and data recovery to defend against physical adversarial attacks, but this approach may not be effective against attacks that use complex perturbations. The limitations of input validation techniques highlight the need for more advanced defense strategies.

Detection-based approaches, such as anomaly detection and classification, can also be used to defend against physical adversarial attacks. These methods aim to detect adversarial examples by identifying patterns or anomalies in the input data that are indicative of an attack. For instance, \cite{doi:10.1145/3450267.3450535} proposed a real-time detector for digital and physical adversarial inputs, but the effectiveness of this approach in defending against physical attacks is still limited. The development of more effective detection-based approaches is crucial for improving the robustness of computer vision systems.

In addition to these defense strategies, various techniques have been proposed to improve the robustness of deep learning models to physical adversarial attacks. For example, \cite{doi:10.1609/aaai.v35i12.17259} proposed a feature space adversarial attack that perturbs the style of the input data, rather than the pixel values. This approach can be used to generate adversarial examples that are more robust to physical transformations, such as changes in lighting or viewpoint. Other techniques, such as \cite{doi:10.1109/cvpr52688.2022.01487} and \cite{doi:10.1109/vr51125.2022.00073}, have been proposed to generate physical adversarial examples that can evade detection by deep learning models.

The development of more effective defense strategies is crucial for ensuring the security and reliability of computer vision systems in the physical world. As \cite{doi:10.1109/iros51168.2021.9636638} and \cite{doi:10.1109/access.2020.2993304} have shown, the development of more robust and effective defense strategies is essential for mitigating the risks associated with visual adversarial attacks. The emerging trends and future directions in this domain will be discussed in the following section, which will delve into the development of more sophisticated attack methods and robust defense strategies. By combining multiple defense strategies and techniques, we can improve the robustness of computer vision systems and ensure their safe and reliable operation in the physical world.


\subsection{Emerging Trends and Future Directions in Visual Adversarial Attacks}


The realm of visual adversarial attacks has witnessed significant growth and evolution, with a plethora of research focused on understanding and mitigating these threats. As we delve into the emerging trends and future directions in this domain, it becomes evident that the development of more sophisticated attack methods and robust defense strategies is crucial \cite{doi:10.1145/3128572.3140444}. One of the key areas of focus is the creation of physical adversarial examples that can deceive machine learning models in real-world scenarios, such as autonomous vehicles and face recognition systems \cite{arxiv:1707.07397}. These examples have the potential to cause significant harm, and therefore, it is essential to develop effective defense mechanisms to counter them.

Recent studies have shown that adversarial training and robust optimization techniques can be used to improve the robustness of machine learning models against visual adversarial attacks \cite{arxiv:1705.07204}. However, these methods are not foolproof, and there is a need for more advanced techniques that can detect and prevent such attacks \cite{arxiv:1902.06705}. The use of input validation and preprocessing techniques has also been explored, but these methods have their limitations and can be evaded by sophisticated attackers \cite{doi:10.1109/tnnls.2021.3105238}.

Another area of research is focused on the development of more realistic and comprehensive evaluation metrics for assessing the robustness of machine learning models against visual adversarial attacks \cite{doi:10.1007/978-3-030-01258-8_39}. This includes the use of metrics such as distortion, success rate, and transferability of adversarial examples between different models \cite{doi:10.1145/3274895.3274904}. The development of standardized benchmarking protocols is also crucial for comparing the robustness of different models and defense strategies \cite{doi:10.1109/tgrs.2022.3225306}.

In addition to these technical advancements, there is a growing need to understand the physical-world implications of visual adversarial attacks. This includes the development of attacks that can be launched in the physical world, such as using stickers or projectors to create adversarial examples \cite{doi:10.1007/978-3-030-58589-1_3}. The use of natural phenomena, such as shadows, to create adversarial examples is also an area of research \cite{arxiv:2203.03818}. These attacks have significant implications for safety-critical systems, such as autonomous vehicles, and highlight the need for more robust defense mechanisms.

The future of visual adversarial attacks is likely to be shaped by the development of more advanced attack methods and defense strategies. The use of emerging technologies, such as edge AI and explainable AI, is expected to play a significant role in this domain \cite{doi:10.24963/ijcai.2021/694}. The development of more realistic and comprehensive evaluation metrics, as well as standardized benchmarking protocols, will also be crucial for advancing the field \cite{arxiv:1902.06705}. As we move forward, it is essential to prioritize the development of robust defense mechanisms that can counter the evolving threats of visual adversarial attacks \cite{arxiv:1707.08945}.

In conclusion, the domain of visual adversarial attacks is rapidly evolving, with significant implications for safety-critical systems. The development of more sophisticated attack methods and robust defense strategies is crucial, and emerging technologies are expected to play a significant role in shaping the future of this domain \cite{doi:10.1109/cvpr42600.2020.01426}. As researchers, it is essential to prioritize the development of robust defense mechanisms and to continue advancing the field through rigorous research and evaluation \cite{arxiv:2209.14262}. By doing so, we can ensure the safe and reliable deployment of machine learning models in a wide range of applications \cite{arxiv:1904.00759}.


\section{Evaluation Metrics and Benchmarking for Adversarial Robustness}



\subsection{Introduction to Evaluation Metrics for Adversarial Robustness}


Evaluating the robustness of deep learning models against adversarial attacks is crucial for ensuring the reliability and security of various applications, including those in the physical world \cite{arxiv:1810.00069}. Existing evaluation metrics, such as accuracy and robustness, have their strengths and limitations. For instance, accuracy metrics may not fully capture the model's ability to withstand adversarial attacks, while robustness metrics may not account for the complexity of real-world scenarios.

In the context of physical-world adversarial attacks, evaluation metrics should consider factors such as environmental conditions, object geometry, and sensor characteristics \cite{arxiv:1707.08945}. The Expectation Over Transformation (EOT) technique has been proposed to enhance the robustness of adversarial perturbations in physical-world scenarios \cite{doi:10.1109/cvpr42600.2020.01426}. 

Recent studies have proposed various evaluation metrics for adversarial robustness, including the use of adversarial training and robust optimization techniques. These metrics aim to improve the model's ability to withstand adversarial attacks by incorporating adversarial examples into the training process \cite{doi:10.24963/ijcai.2021/694}. 

The development of standardized benchmarking protocols is also essential for comparing the robustness of different models and defense strategies \cite{doi:10.1109/tgrs.2022.3225306}. Existing benchmarking protocols, such as the APRICOT dataset, provide a comprehensive evaluation framework for physical-world adversarial attacks \cite{doi:10.1007/978-3-030-58589-1_3}.

In conclusion, evaluating the robustness of deep learning models against adversarial attacks in physical-world scenarios requires comprehensive and realistic evaluation metrics \cite{doi:10.1109/access.2022.3208131}. Future research directions include the development of more realistic and comprehensive evaluation metrics, the incorporation of human perception and environmental factors, and the establishment of standardized benchmarking protocols for comparing the robustness of different models and defense strategies \cite{doi:10.1007/s11633-019-1211-x}.


\subsection{Benchmarking Protocols for Adversarial Robustness}


The development of standardized benchmarking protocols is crucial for evaluating and comparing the robustness of different models and defense strategies against adversarial attacks in physical-world scenarios. As highlighted in \cite{arxiv:1707.07397}, the existence of robust 3D adversarial objects poses a significant threat to real-world systems, emphasizing the need for comprehensive benchmarking protocols that can accurately assess the effectiveness of various defense strategies. To address this challenge, researchers have proposed the use of evaluation metrics specifically designed for physical-world adversarial attacks, such as success rate, distance metric (e.g., L2, L$\infty$), and perceptual path length, as mentioned in \cite{doi:10.24963/ijcai.2021/694}. These metrics can help evaluate the effectiveness and robustness of physical-world adversarial attacks, which is essential for developing more robust and reliable deep learning models.

Existing benchmarking protocols, such as those proposed in \cite{doi:10.1109/access.2018.2807385}, have limitations in evaluating adversarial robustness in physical-world scenarios. For instance, they often focus on digital attacks and do not account for the complexities of real-world environments, including environmental conditions, object geometry, and sensor characteristics. To overcome these limitations, researchers have proposed the use of techniques such as expectation over transformation (EOT), as discussed in \cite{doi:10.1007/978-3-030-10925-7_4}, which can help improve the robustness of physical adversarial examples to various transformations and perturbations.

The development of benchmarking protocols for physical-world adversarial attacks requires careful consideration of various factors, including environmental conditions, object geometry, and sensor characteristics. As noted in \cite{doi:10.1109/lcomm.2019.2901469}, the openness of the wireless channel and the presence of noise and interference can significantly impact the effectiveness of adversarial attacks. Furthermore, the selection of appropriate evaluation metrics is critical for assessing the robustness of models against physical-world adversarial attacks. As highlighted in \cite{doi:10.1126/science.aaw4399}, the use of metrics such as accuracy, robustness, and attack success rate can provide valuable insights into the effectiveness of defense strategies.

In addition to the design of benchmarking protocols, the development of more comprehensive and realistic evaluation metrics is essential for accurately assessing the robustness of models in physical-world scenarios. As discussed in \cite{doi:10.1049/cit2.12028}, the development of evaluation metrics that incorporate human perception and environmental factors is crucial for evaluating the effectiveness of physical-world adversarial attacks. The use of advanced techniques, such as generative adversarial networks (GANs) and reinforcement learning, can help improve the robustness and effectiveness of physical adversarial examples, as discussed in \cite{doi:10.1109/cvpr46437.2021.00846}.

In conclusion, the development of standardized benchmarking protocols for physical-world adversarial attacks is essential for evaluating and comparing the robustness of different models and defense strategies. The design of such protocols requires careful consideration of various factors, including environmental conditions, object geometry, and sensor characteristics. The selection of appropriate evaluation metrics, such as success rate, distance metric, and perceptual path length, is critical for assessing the robustness of models against physical-world adversarial attacks. As highlighted in \cite{doi:10.1007/978-3-030-58589-1_3}, the development of more comprehensive and realistic evaluation metrics, such as those that incorporate human perception and environmental factors, is essential for accurately assessing the robustness of models in physical-world scenarios. Future research should focus on addressing the challenges and limitations of existing benchmarking protocols and developing more effective and realistic evaluation metrics for physical-world adversarial attacks, as discussed in \cite{doi:10.1016/j.eng.2019.12.012}.


\subsection{Evaluation Metrics for Physical-World Adversarial Attacks}


The development of evaluation metrics specifically designed for physical-world adversarial attacks is a crucial aspect of assessing the effectiveness and robustness of such attacks. As \cite{doi:10.24963/ijcai.2021/694} notes, the existence of adversarial examples in the physical world poses significant security concerns, and evaluating their impact requires tailored metrics. In this context, metrics such as success rate, distance metric (e.g., L2, L$\infty$), and perceptual path length are essential in capturing the effectiveness, stealthiness, and robustness of physical-world adversarial attacks \cite{arxiv:1707.07397}.

One key challenge in evaluating physical-world adversarial attacks is the need to consider human perception. As \cite{doi:10.1145/3128572.3140444} highlights, adversarial examples can be designed to be imperceptible to humans, making it essential to develop metrics that account for human visual perception. To address this, researchers have proposed metrics such as the structural similarity index (SSIM) and the Frechet inception distance (FID), which can help evaluate the perceptual quality of adversarial examples \cite{doi:10.1109/iccvw.2019.00257}.

Another critical aspect of evaluating physical-world adversarial attacks is the consideration of environmental factors. As \cite{arxiv:1707.08945} notes, physical-world attacks must be robust to various environmental conditions, such as lighting, viewpoint, and weather. To capture this, metrics such as the expectation over transformation (EOT) can be used to evaluate the robustness of adversarial examples under different environmental conditions \cite{doi:10.1007/978-3-030-10925-7_4}.

In addition to these metrics, researchers have also proposed evaluation frameworks that consider multiple factors, such as attack success rate, perturbation size, and robustness to transformations \cite{arxiv:1802.06430}. These frameworks provide a comprehensive approach to evaluating physical-world adversarial attacks and can help identify the most effective and robust attacks.

The development of evaluation metrics for physical-world adversarial attacks is an active area of research, with new metrics and frameworks being proposed regularly. As \cite{doi:10.1186/s42400-018-0012-9} notes, the evaluation of physical-world adversarial attacks is a complex task that requires careful consideration of multiple factors. 

Furthermore, the evaluation of physical-world adversarial attacks should also consider the potential defenses against such attacks. As \cite{doi:10.1145/3274895.3274904} highlights, the development of effective defenses against physical-world adversarial attacks is crucial for ensuring the security and reliability of deep learning systems. Researchers have proposed various defense methods, including adversarial training, input validation, and detection-based approaches \cite{arxiv:1910.11099}.

In conclusion, the development of evaluation metrics for physical-world adversarial attacks is a critical aspect of assessing the effectiveness and robustness of such attacks. By considering human perception, environmental factors, and potential defenses, researchers can develop comprehensive and realistic evaluation metrics that can help identify the most effective and robust attacks. As \cite{doi:10.1109/cvpr42600.2020.01426} notes, the evaluation of physical-world adversarial attacks is essential for ensuring the security and reliability of deep learning systems in safety-critical applications. Future research should continue to focus on developing more advanced evaluation metrics and frameworks that can capture the complexities of physical-world adversarial attacks.


\subsection{Adversarial Robustness Assessment in Real-World Scenarios}


Assessing adversarial robustness in real-world scenarios is a complex task that requires careful consideration of various factors, including the type of attack, the environment, and the model's architecture. As \cite{doi:10.1201/9781351251389-8} have shown, even small perturbations in the physical world can cause significant errors in deep learning models. To evaluate the robustness of models in real-world scenarios, researchers have proposed various methodologies, including ensemble methods, transferability attacks, and uncertainty estimation. These methodologies are crucial in identifying vulnerabilities in individual models and providing a more comprehensive understanding of the robustness of the overall system.

Ensemble methods involve combining multiple models to improve robustness, as demonstrated in \cite{arxiv:2110.03825}. This approach can help to identify vulnerabilities in individual models and provide a more comprehensive understanding of the robustness of the overall system. For instance, ensemble methods can be used to evaluate the robustness of models to different types of attacks, such as physical-world attacks or digital attacks. Transferability attacks, on the other hand, involve attacking a model with adversarial examples crafted for a different model, as shown in \cite{doi:10.1109/iros51168.2021.9636638}. This approach can help to evaluate the robustness of models to attacks that are not specifically designed for them, which is an important consideration in real-world scenarios.

Uncertainty estimation is another approach that has been proposed to evaluate adversarial robustness, as discussed in \cite{doi:10.1109/access.2020.2993304}. This approach involves estimating the uncertainty of a model's predictions and using this information to detect potential adversarial examples. Uncertainty estimation can be particularly useful in real-world scenarios where the model is faced with uncertain or noisy inputs. By combining uncertainty estimation with ensemble methods and transferability attacks, researchers can gain a more comprehensive understanding of the robustness of models in real-world scenarios.

In addition to these methodologies, researchers have also proposed various metrics to evaluate adversarial robustness, including the adversarial loss, the robustness metric, and the uncertainty metric, as discussed in \cite{doi:10.1109/wacv51458.2022.00144}. These metrics can provide a quantitative evaluation of a model's robustness to adversarial attacks and can be used to compare the robustness of different models. However, the development of these metrics is closely tied to the development of evaluation protocols, which is a critical aspect of assessing adversarial robustness in real-world scenarios.

Despite the progress made in evaluating adversarial robustness, there are still many challenges that need to be addressed. One of the main challenges is the lack of standardized evaluation protocols, as noted in \cite{doi:10.1007/978-3-030-58589-1_3}. This makes it difficult to compare the robustness of different models and to evaluate the effectiveness of different defense mechanisms. Another challenge is the need for more realistic and comprehensive evaluation metrics that can capture the complexity of real-world scenarios, as discussed in \cite{doi:10.1609/aaai.v35i2.16211}. To address these challenges, researchers have proposed various future directions, including the development of more realistic and comprehensive evaluation metrics, the creation of standardized evaluation protocols, and the investigation of new defense mechanisms, as discussed in \cite{arxiv:2007.04137}.

The investigation of new defense mechanisms is crucial in improving adversarial robustness, and emerging technologies such as edge AI and explainable AI may provide new opportunities for improving adversarial robustness, as noted in \cite{doi:10.1109/asp-dac47756.2020.9045584}. Additionally, the use of physical-world attacks, such as those discussed in \cite{doi:10.1109/lra.2022.3189783}, \cite{arxiv:2210.15291}, and \cite{doi:10.1109/vr51125.2022.00073}, can help to evaluate the robustness of models in real-world scenarios. By addressing the challenges and limitations of current evaluation methodologies and metrics, researchers can develop more effective and robust models that can withstand adversarial attacks in the physical world. This will be further explored in the following sections, which will discuss the development of evaluation metrics and benchmarking protocols for adversarial robustness.


\subsection{Future Directions in Evaluation Metrics and Benchmarking}


The development of evaluation metrics and benchmarking protocols for adversarial robustness is an active area of research, with significant implications for the security and reliability of deep learning systems. As \cite{doi:10.1145/3128572.3140444} and \cite{arxiv:1902.06705} highlight, existing metrics and protocols have limitations, and there is a need for more comprehensive and realistic evaluations. 

Another promising area of research is the development of metrics that capture the robustness of models to different types of perturbations, such as \cite{arxiv:1801.02612}. These metrics can help evaluate the effectiveness of defense mechanisms and identify potential vulnerabilities in deep learning systems. Furthermore, \cite{doi:10.1609/aaai.v35i9.16927} shows that attribute-guided adversarial training can improve the robustness of models to natural perturbations.

The use of benchmarking protocols, such as \cite{doi:10.1109/tgrs.2022.3225306}, can facilitate the comparison of different defense mechanisms and evaluation metrics. These protocols can help identify the most effective defense strategies and provide a comprehensive understanding of the robustness of deep learning systems. Additionally, \cite{arxiv:1707.08945} demonstrates the importance of considering physical-world attacks in the evaluation of adversarial robustness.

Emerging trends in adversarial robustness evaluation include the use of \cite{arxiv:1904.00759} and \cite{arxiv:2203.03818}. These attacks highlight the need for more comprehensive evaluation metrics and protocols that consider the physical world and natural phenomena. 

In conclusion, the development of evaluation metrics and benchmarking protocols for adversarial robustness is a critical area of research, with significant implications for the security and reliability of deep learning systems. Future research directions include the incorporation of emerging technologies, the development of more comprehensive and realistic metrics, and the consideration of physical-world attacks. By synthesizing information from various studies, including \cite{doi:10.24963/ijcai.2021/694} and \cite{doi:10.1007/978-3-030-58589-1_3}, we can gain a deeper understanding of the challenges and opportunities in this field, and work towards developing more robust and reliable deep learning systems that can withstand adversarial attacks in the physical world. Ultimately, the goal of this research is to provide a foundation for the development of more comprehensive and realistic evaluation metrics and benchmarking protocols, and to facilitate the creation of more robust and reliable deep learning systems.


\section{Conclusion}


The field of visual adversarial attacks and defenses in the physical world has witnessed significant advancements in recent years, with a plethora of research focusing on the vulnerabilities of deep learning models to carefully crafted inputs. As evident from \cite{doi:10.24963/ijcai.2021/694}, \cite{arxiv:1707.08945}, and \cite{doi:10.1109/access.2018.2807385}, the security concerns associated with these attacks have become a pressing issue, particularly in safety-critical applications such as autonomous vehicles and face recognition systems. The development of robust defense mechanisms against such attacks is crucial to ensure the reliability and trustworthiness of deep learning-based systems.

Comparative analysis of different approaches reveals that adversarial training and robust optimization techniques are effective in improving the robustness of deep learning models against adversarial attacks. However, these methods often come with a trade-off in terms of computational cost and accuracy. Ensemble methods and hybrid approaches have also shown promise in defending against visual adversarial attacks. Nevertheless, the evaluation metrics and benchmarking protocols for assessing adversarial robustness are still in their infancy, and more research is needed to develop comprehensive and realistic evaluation metrics.

Emerging trends in the field include the use of edge AI and explainable AI to enhance adversarial robustness. The application of physical adversarial attacks to other domains, such as audio and graph data, is also an area of increasing interest. Furthermore, the development of novel attack methods, such as \cite{doi:10.1109/vr51125.2022.00073} and \cite{arxiv:2203.08392}, highlights the need for continued research into the vulnerabilities of deep learning models.

In conclusion, the field of visual adversarial attacks and defenses in the physical world is rapidly evolving, with new challenges and opportunities emerging continuously. As noted in \cite{doi:10.1145/3433210.3437513} and \cite{doi:10.1109/jiot.2021.3099164}, the development of robust defense mechanisms against adversarial attacks is crucial to ensure the security and reliability of deep learning-based systems. Future research directions should focus on developing comprehensive evaluation metrics, improving the robustness of deep learning models, and exploring the application of physical adversarial attacks to other domains. By synthesizing the insights from existing research, such as \cite{doi:10.1109/iros51168.2021.9636638} and \cite{doi:10.1109/cvpr42600.2020.01426}, we can work towards creating more robust and trustworthy deep learning-based systems. Ultimately, the development of effective defense mechanisms against visual adversarial attacks will require a multidisciplinary approach, combining advances in deep learning, computer vision, and cybersecurity \cite{doi:10.1109/cvpr42600.2020.00108}.

\begin{thebibliography}{91}
\setlength{\itemsep}{2pt}

\bibitem{doi:10.1201/9781351251389-8}
\newblock \emph{Adversarial examples in the physical world}.
\newblock DOI:10.1201/9781351251389-8, 2018.
\newblock \url{https://doi.org/10.1201/9781351251389-8}

\bibitem{arxiv:1707.08945}
\newblock \emph{Robust Physical-World Attacks on Deep Learning Models}.
\newblock arXiv:1707.08945, 2017.
\newblock \url{https://arxiv.org/abs/1707.08945}

\bibitem{doi:10.24963/ijcai.2021/694}
\newblock \emph{Adversarial Examples in Physical World}.
\newblock DOI:10.24963/ijcai.2021/694, 2021.
\newblock \url{https://doi.org/10.24963/ijcai.2021/694}

\bibitem{doi:10.1109/cvpr42600.2020.01426}
\newblock \emph{PhysGAN: Generating Physical-World-Resilient Adversarial Examples for Autonomous Driving}.
\newblock DOI:10.1109/cvpr42600.2020.01426, 2020.
\newblock \url{https://doi.org/10.1109/cvpr42600.2020.01426}

\bibitem{doi:10.1016/j.sysarc.2020.101766}
\newblock \emph{Attacking Vision-based Perception in End-to-End Autonomous Driving Models}.
\newblock DOI:10.1016/j.sysarc.2020.101766, 2020.
\newblock \url{https://doi.org/10.1016/j.sysarc.2020.101766}

\bibitem{doi:10.1109/jiot.2021.3099164}
\newblock \emph{Evaluating Adversarial Attacks on Driving Safety in Vision-Based Autonomous Vehicles}.
\newblock DOI:10.1109/jiot.2021.3099164, 2021.
\newblock \url{https://doi.org/10.1109/jiot.2021.3099164}

\bibitem{doi:10.1016/j.eng.2019.12.012}
\newblock \emph{Adversarial Attacks and Defenses in Deep Learning}.
\newblock DOI:10.1016/j.eng.2019.12.012, 2020.
\newblock \url{https://doi.org/10.1016/j.eng.2019.12.012}

\bibitem{arxiv:2102.01356}
\newblock \emph{Recent Advances in Adversarial Training for Adversarial Robustness}.
\newblock arXiv:2102.01356, 2021.
\newblock \url{https://arxiv.org/abs/2102.01356}

\bibitem{doi:10.1109/iros51168.2021.9636638}
\newblock \emph{Adversarial Attacks on Camera-LiDAR Models for 3D Car Detection}.
\newblock DOI:10.1109/iros51168.2021.9636638, 2021.
\newblock \url{https://doi.org/10.1109/iros51168.2021.9636638}

\bibitem{arxiv:1906.11897}
\newblock \emph{On Physical Adversarial Patches for Object Detection}.
\newblock arXiv:1906.11897, 2019.
\newblock \url{https://arxiv.org/abs/1906.11897}

\bibitem{doi:10.1609/aaai.v32i1.12341}
\newblock \emph{Unravelling Robustness of Deep Learning based Face Recognition Against Adversarial Attacks}.
\newblock DOI:10.1609/aaai.v32i1.12341, 2018.
\newblock \url{https://doi.org/10.1609/aaai.v32i1.12341}

\bibitem{arxiv:2007.04137}
\newblock \emph{SLAP: Improving Physical Adversarial Examples with Short-Lived Adversarial Perturbations}.
\newblock arXiv:2007.04137, 2020.
\newblock \url{https://arxiv.org/abs/2007.04137}

\bibitem{doi:10.1109/access.2022.3208131}
\newblock \emph{Adversarial Deep Learning: A Survey on Adversarial Attacks and Defense Mechanisms on Image Classification}.
\newblock DOI:10.1109/access.2022.3208131, 2022.
\newblock \url{https://doi.org/10.1109/access.2022.3208131}

\bibitem{doi:10.1109/asp-dac47756.2020.9045584}
\newblock \emph{LanCe: A Comprehensive and Lightweight CNN Defense Methodology against Physical Adversarial Attacks on Embedded Multimedia Applications}.
\newblock DOI:10.1109/asp-dac47756.2020.9045584, 2020.
\newblock \url{https://doi.org/10.1109/asp-dac47756.2020.9045584}

\bibitem{doi:10.1007/s11633-019-1211-x}
\newblock \emph{Adversarial Attacks and Defenses in Images, Graphs and Text: A Review}.
\newblock DOI:10.1007/s11633-019-1211-x, 2020.
\newblock \url{https://doi.org/10.1007/s11633-019-1211-x}

\bibitem{doi:10.1145/3459637.3482029}
\newblock \emph{Adversarial Robustness of Deep Learning: Theory, Algorithms, and Applications}.
\newblock DOI:10.1145/3459637.3482029, 2021.
\newblock \url{https://doi.org/10.1145/3459637.3482029}

\bibitem{arxiv:1707.07397}
\newblock \emph{Synthesizing Robust Adversarial Examples}.
\newblock arXiv:1707.07397, 2017.
\newblock \url{https://arxiv.org/abs/1707.07397}

\bibitem{doi:10.1109/access.2018.2807385}
\newblock \emph{Threat of Adversarial Attacks on Deep Learning in Computer Vision: A Survey}.
\newblock DOI:10.1109/access.2018.2807385, 2018.
\newblock \url{https://doi.org/10.1109/access.2018.2807385}

\bibitem{doi:10.1109/cvpr42600.2020.00108}
\newblock \emph{Adversarial Camouflage: Hiding Physical-World Attacks With Natural Styles}.
\newblock DOI:10.1109/cvpr42600.2020.00108, 2020.
\newblock \url{https://doi.org/10.1109/cvpr42600.2020.00108}

\bibitem{arxiv:1904.00759}
\newblock \emph{Adversarial camera stickers: A physical camera-based attack on deep learning systems}.
\newblock arXiv:1904.00759, 2019.
\newblock \url{https://arxiv.org/abs/1904.00759}

\bibitem{doi:10.1007/978-3-030-10925-7_4}
\newblock \emph{ShapeShifter: Robust Physical Adversarial Attack on Faster R-CNN Object Detector}.
\newblock DOI:10.1007/978-3-030-10925-7\_4, 2019.
\newblock \url{https://doi.org/10.1007/978-3-030-10925-7_4}

\bibitem{arxiv:1707.03501}
\newblock \emph{NO Need to Worry about Adversarial Examples in Object Detection in Autonomous Vehicles}.
\newblock arXiv:1707.03501, 2017.
\newblock \url{https://arxiv.org/abs/1707.03501}

\bibitem{doi:10.1109/lcomm.2019.2901469}
\newblock \emph{Physical Adversarial Attacks Against End-to-End Autoencoder Communication Systems}.
\newblock DOI:10.1109/lcomm.2019.2901469, 2019.
\newblock \url{https://doi.org/10.1109/lcomm.2019.2901469}

\bibitem{doi:10.1049/cit2.12028}
\newblock \emph{A survey on adversarial attacks and defences}.
\newblock DOI:10.1049/cit2.12028, 2021.
\newblock \url{https://doi.org/10.1049/cit2.12028}

\bibitem{doi:10.1109/iccvw.2019.00257}
\newblock \emph{AdvGAN++: Harnessing Latent Layers for Adversary Generation}.
\newblock DOI:10.1109/iccvw.2019.00257, 2019.
\newblock \url{https://doi.org/10.1109/iccvw.2019.00257}

\bibitem{arxiv:1807.07769}
\newblock \emph{Physical Adversarial Examples for Object Detectors}.
\newblock arXiv:1807.07769, 2018.
\newblock \url{https://arxiv.org/abs/1807.07769}

\bibitem{doi:10.1609/aaai.v34i01.5443}
\newblock \emph{Robust Adversarial Objects against Deep Learning Models}.
\newblock DOI:10.1609/aaai.v34i01.5443, 2020.
\newblock \url{https://doi.org/10.1609/aaai.v34i01.5443}

\bibitem{arxiv:1802.06430}
\newblock \emph{DARTS: Deceiving Autonomous Cars with Toxic Signs}.
\newblock arXiv:1802.06430, 2018.
\newblock \url{https://arxiv.org/abs/1802.06430}

\bibitem{doi:10.1109/biosig52210.2021.9548291}
\newblock \emph{On Brightness Agnostic Adversarial Examples Against Face Recognition Systems}.
\newblock DOI:10.1109/biosig52210.2021.9548291, 2021.
\newblock \url{https://doi.org/10.1109/biosig52210.2021.9548291}

\bibitem{doi:10.1007/978-3-030-58589-1_3}
\newblock \emph{APRICOT: A Dataset of Physical Adversarial Attacks on Object Detection}.
\newblock DOI:10.1007/978-3-030-58589-1\_3, 2020.
\newblock \url{https://doi.org/10.1007/978-3-030-58589-1_3}

\bibitem{arxiv:2001.05574}
\newblock \emph{Advbox: a toolbox to generate adversarial examples that fool neural networks}.
\newblock arXiv:2001.05574, 2020.
\newblock \url{https://arxiv.org/abs/2001.05574}

\bibitem{doi:10.1109/tip.2021.3092822}
\newblock \emph{Toward Visual Distortion in Black-Box Attacks}.
\newblock DOI:10.1109/tip.2021.3092822, 2021.
\newblock \url{https://doi.org/10.1109/tip.2021.3092822}

\bibitem{doi:10.1109/cvpr42600.2020.00049}
\newblock \emph{One Man's Trash Is Another Man's Treasure: Resisting Adversarial Examples by Adversarial Examples}.
\newblock DOI:10.1109/cvpr42600.2020.00049, 2020.
\newblock \url{https://doi.org/10.1109/cvpr42600.2020.00049}

\bibitem{doi:10.1016/j.neucom.2022.05.031}
\newblock \emph{Enhancing real-world adversarial patches through 3D modeling of complex target scenes}.
\newblock DOI:10.1016/j.neucom.2022.05.031, 2022.
\newblock \url{https://doi.org/10.1016/j.neucom.2022.05.031}

\bibitem{arxiv:1910.11099}
\newblock \emph{Adversarial T-shirt! Evading Person Detectors in A Physical World}.
\newblock arXiv:1910.11099, 2019.
\newblock \url{https://arxiv.org/abs/1910.11099}

\bibitem{doi:10.1109/cvpr52688.2022.01487}
\newblock \emph{DTA: Physical Camouflage Attacks using Differentiable Transformation Network}.
\newblock DOI:10.1109/cvpr52688.2022.01487, 2022.
\newblock \url{https://doi.org/10.1109/cvpr52688.2022.01487}

\bibitem{doi:10.1109/cvpr46437.2021.00614}
\newblock \emph{Backdoor Attacks Against Deep Learning Systems in the Physical World}.
\newblock DOI:10.1109/cvpr46437.2021.00614, 2021.
\newblock \url{https://doi.org/10.1109/cvpr46437.2021.00614}

\bibitem{doi:10.1145/3319535.3339815}
\newblock \emph{Adversarial Sensor Attack on LiDAR-based Perception in Autonomous Driving}.
\newblock DOI:10.1145/3319535.3339815, 2019.
\newblock \url{https://doi.org/10.1145/3319535.3339815}

\bibitem{arxiv:2007.16118}
\newblock \emph{Physical Adversarial Attack on Vehicle Detector in the Carla Simulator}.
\newblock arXiv:2007.16118, 2020.
\newblock \url{https://arxiv.org/abs/2007.16118}

\bibitem{doi:10.1145/3128572.3140444}
\newblock \emph{Adversarial Examples Are Not Easily Detected: Bypassing Ten Detection Methods}.
\newblock DOI:10.1145/3128572.3140444, 2017.
\newblock \url{https://doi.org/10.1145/3128572.3140444}

\bibitem{arxiv:1904.02884}
\newblock \emph{Evading Defenses to Transferable Adversarial Examples by Translation-Invariant Attacks}.
\newblock arXiv:1904.02884, 2019.
\newblock \url{https://arxiv.org/abs/1904.02884}

\bibitem{doi:10.1016/j.patcog.2022.109009}
\newblock \emph{Robust Physical-World Attacks on Face Recognition}.
\newblock DOI:10.1016/j.patcog.2022.109009, 2022.
\newblock \url{https://doi.org/10.1016/j.patcog.2022.109009}

\bibitem{arxiv:1904.13000}
\newblock \emph{Adversarial Training and Robustness for Multiple Perturbations}.
\newblock arXiv:1904.13000, 2019.
\newblock \url{https://arxiv.org/abs/1904.13000}

\bibitem{doi:10.1609/aaai.v32i1.11828}
\newblock \emph{Detecting Adversarial Examples Through Image Transformation}.
\newblock DOI:10.1609/aaai.v32i1.11828, 2018.
\newblock \url{https://doi.org/10.1609/aaai.v32i1.11828}

\bibitem{doi:10.1109/tnnls.2021.3105238}
\newblock \emph{Detecting Adversarial Examples by Input Transformations, Defense Perturbations, and Voting}.
\newblock DOI:10.1109/tnnls.2021.3105238, 2021.
\newblock \url{https://doi.org/10.1109/tnnls.2021.3105238}

\bibitem{arxiv:1902.06705}
\newblock \emph{On Evaluating Adversarial Robustness}.
\newblock arXiv:1902.06705, 2019.
\newblock \url{https://arxiv.org/abs/1902.06705}

\bibitem{arxiv:2203.03818}
\newblock \emph{Shadows can be Dangerous: Stealthy and Effective Physical-world Adversarial Attack by Natural Phenomenon}.
\newblock arXiv:2203.03818, 2022.
\newblock \url{https://arxiv.org/abs/2203.03818}

\bibitem{arxiv:2209.14262}
\newblock \emph{A Survey on Physical Adversarial Attack in Computer Vision}.
\newblock arXiv:2209.14262, 2022.
\newblock \url{https://arxiv.org/abs/2209.14262}

\bibitem{doi:10.1016/j.jisa.2020.102694}
\newblock \emph{Robust and Natural Physical Adversarial Examples for Object Detectors}.
\newblock DOI:10.1016/j.jisa.2020.102694, 2021.
\newblock \url{https://doi.org/10.1016/j.jisa.2020.102694}

\bibitem{arxiv:1712.07107}
\newblock \emph{Adversarial Examples: Attacks and Defenses for Deep Learning}.
\newblock arXiv:1712.07107, 2017.
\newblock \url{https://arxiv.org/abs/1712.07107}

\bibitem{arxiv:1712.09665}
\newblock \emph{Adversarial Patch}.
\newblock arXiv:1712.09665, 2017.
\newblock \url{https://arxiv.org/abs/1712.09665}

\bibitem{arxiv:1911.05268}
\newblock \emph{Adversarial Examples in Modern Machine Learning: A Review}.
\newblock arXiv:1911.05268, 2019.
\newblock \url{https://arxiv.org/abs/1911.05268}

\bibitem{doi:10.1109/tgrs.2022.3225306}
\newblock \emph{Benchmarking Adversarial Patch Against Aerial Detection}.
\newblock DOI:10.1109/tgrs.2022.3225306, 2022.
\newblock \url{https://doi.org/10.1109/tgrs.2022.3225306}

\bibitem{arxiv:2104.02361}
\newblock \emph{Backdoor Attack in the Physical World}.
\newblock arXiv:2104.02361, 2021.
\newblock \url{https://arxiv.org/abs/2104.02361}

\bibitem{arxiv:2209.09502}
\newblock \emph{GAMA: Generative Adversarial Multi-Object Scene Attacks}.
\newblock arXiv:2209.09502, 2022.
\newblock \url{https://arxiv.org/abs/2209.09502}

\bibitem{arxiv:2101.07922}
\newblock \emph{LowKey: Leveraging Adversarial Attacks to Protect Social Media Users from Facial Recognition}.
\newblock arXiv:2101.07922, 2021.
\newblock \url{https://arxiv.org/abs/2101.07922}

\bibitem{doi:10.1145/3433210.3437513}
\newblock \emph{ConAML: Constrained Adversarial Machine Learning for Cyber-Physical Systems}.
\newblock DOI:10.1145/3433210.3437513, 2021.
\newblock \url{https://doi.org/10.1145/3433210.3437513}

\bibitem{arxiv:1810.00069}
\newblock \emph{Adversarial Attacks and Defences: A Survey}.
\newblock arXiv:1810.00069, 2018.
\newblock \url{https://arxiv.org/abs/1810.00069}

\bibitem{doi:10.1109/cvpr.2019.00790}
\newblock \emph{Efficient Decision-Based Black-Box Adversarial Attacks on Face Recognition}.
\newblock DOI:10.1109/cvpr.2019.00790, 2019.
\newblock \url{https://doi.org/10.1109/cvpr.2019.00790}

\bibitem{doi:10.1145/3377811.3380422}
\newblock \emph{DeepBillboard: Systematic Physical-World Testing of Autonomous Driving Systems}.
\newblock DOI:10.1145/3377811.3380422, 2020.
\newblock \url{https://doi.org/10.1145/3377811.3380422}

\bibitem{doi:10.1145/3274895.3274904}
\newblock \emph{Adversarial examples in remote sensing}.
\newblock DOI:10.1145/3274895.3274904, 2018.
\newblock \url{https://doi.org/10.1145/3274895.3274904}

\bibitem{doi:10.1145/3398394}
\newblock \emph{Adversarial Examples on Object Recognition}.
\newblock DOI:10.1145/3398394, 2020.
\newblock \url{https://doi.org/10.1145/3398394}

\bibitem{doi:10.1080/00031305.2021.2006781}
\newblock \emph{A Review of Adversarial Attack and Defense for Classification Methods}.
\newblock DOI:10.1080/00031305.2021.2006781, 2021.
\newblock \url{https://doi.org/10.1080/00031305.2021.2006781}

\bibitem{doi:10.1109/access.2022.3160283}
\newblock \emph{ARGAN: Adversarially Robust Generative Adversarial Networks for Deep Neural Networks against Adversarial Examples}.
\newblock DOI:10.1109/access.2022.3160283, 2022.
\newblock \url{https://doi.org/10.1109/access.2022.3160283}

\bibitem{doi:10.1109/vr51125.2022.00073}
\newblock \emph{SPAA: Stealthy Projector-based Adversarial Attacks on Deep Image Classifiers}.
\newblock DOI:10.1109/vr51125.2022.00073, 2022.
\newblock \url{https://doi.org/10.1109/vr51125.2022.00073}

\bibitem{doi:10.1007/s13042-021-01435-0}
\newblock \emph{Really natural adversarial examples}.
\newblock DOI:10.1007/s13042-021-01435-0, 2021.
\newblock \url{https://doi.org/10.1007/s13042-021-01435-0}

\bibitem{doi:10.1145/3450267.3450535}
\newblock \emph{Real-time detectors for digital and physical adversarial inputs to perception systems}.
\newblock DOI:10.1145/3450267.3450535, 2021.
\newblock \url{https://doi.org/10.1145/3450267.3450535}

\bibitem{arxiv:1705.07204}
\newblock \emph{Ensemble Adversarial Training: Attacks and Defenses}.
\newblock arXiv:1705.07204, 2017.
\newblock \url{https://arxiv.org/abs/1705.07204}

\bibitem{arxiv:1901.08573}
\newblock \emph{Theoretically Principled Trade-off between Robustness and Accuracy}.
\newblock arXiv:1901.08573, 2019.
\newblock \url{https://arxiv.org/abs/1901.08573}

\bibitem{arxiv:2012.12235}
\newblock \emph{Unadversarial Examples: Designing Objects for Robust Vision}.
\newblock arXiv:2012.12235, 2020.
\newblock \url{https://arxiv.org/abs/2012.12235}

\bibitem{doi:10.1609/aaai.v35i9.16927}
\newblock \emph{Attribute-Guided Adversarial Training for Robustness to Natural Perturbations}.
\newblock DOI:10.1609/aaai.v35i9.16927, 2021.
\newblock \url{https://doi.org/10.1609/aaai.v35i9.16927}

\bibitem{doi:10.1109/eurosp53844.2022.00047}
\newblock \emph{GRAPHITE: Generating Automatic Physical Examples for Machine-Learning Attacks on Computer Vision Systems}.
\newblock DOI:10.1109/eurosp53844.2022.00047, 2022.
\newblock \url{https://doi.org/10.1109/eurosp53844.2022.00047}

\bibitem{doi:10.1126/science.aaw4399}
\newblock \emph{Adversarial attacks on medical machine learning}.
\newblock DOI:10.1126/science.aaw4399, 2019.
\newblock \url{https://doi.org/10.1126/science.aaw4399}

\bibitem{doi:10.1155/2021/4907754}
\newblock \emph{A Survey on Adversarial Attack in the Age of Artificial Intelligence}.
\newblock DOI:10.1155/2021/4907754, 2021.
\newblock \url{https://doi.org/10.1155/2021/4907754}

\bibitem{arxiv:1711.08478}
\newblock \emph{MagNet and "Efficient Defenses Against Adversarial Attacks" are Not Robust to Adversarial Examples}.
\newblock arXiv:1711.08478, 2017.
\newblock \url{https://arxiv.org/abs/1711.08478}

\bibitem{doi:10.1186/s42400-018-0012-9}
\newblock \emph{A survey of practical adversarial example attacks}.
\newblock DOI:10.1186/s42400-018-0012-9, 2018.
\newblock \url{https://doi.org/10.1186/s42400-018-0012-9}

\bibitem{doi:10.1609/aaai.v35i2.16211}
\newblock \emph{Towards Universal Physical Attacks on Single Object Tracking}.
\newblock DOI:10.1609/aaai.v35i2.16211, 2021.
\newblock \url{https://doi.org/10.1609/aaai.v35i2.16211}

\bibitem{doi:10.1007/978-3-030-58548-8_1}
\newblock \emph{Making an Invisibility Cloak: Real World Adversarial Attacks on Object Detectors}.
\newblock DOI:10.1007/978-3-030-58548-8\_1, 2020.
\newblock \url{https://doi.org/10.1007/978-3-030-58548-8_1}

\bibitem{doi:10.1609/aaai.v35i12.17259}
\newblock \emph{Towards Feature Space Adversarial Attack by Style Perturbation}.
\newblock DOI:10.1609/aaai.v35i12.17259, 2021.
\newblock \url{https://doi.org/10.1609/aaai.v35i12.17259}

\bibitem{arxiv:2110.03825}
\newblock \emph{Exploring Architectural Ingredients of Adversarially Robust Deep Neural Networks}.
\newblock arXiv:2110.03825, 2021.
\newblock \url{https://arxiv.org/abs/2110.03825}

\bibitem{doi:10.1109/wacv51458.2022.00144}
\newblock \emph{Adversarial Robustness of Deep Sensor Fusion Models}.
\newblock DOI:10.1109/wacv51458.2022.00144, 2022.
\newblock \url{https://doi.org/10.1109/wacv51458.2022.00144}

\bibitem{doi:10.1109/ijcnn55064.2022.9892612}
\newblock \emph{fakeWeather: Adversarial Attacks for Deep Neural Networks Emulating Weather Conditions on the Camera Lens of Autonomous Systems}.
\newblock DOI:10.1109/ijcnn55064.2022.9892612, 2022.
\newblock \url{https://doi.org/10.1109/ijcnn55064.2022.9892612}

\bibitem{doi:10.1007/978-3-030-01258-8_39}
\newblock \emph{Is Robustness the Cost of Accuracy? - A Comprehensive Study on the Robustness of 18 Deep Image Classification Models}.
\newblock DOI:10.1007/978-3-030-01258-8\_39, 2018.
\newblock \url{https://doi.org/10.1007/978-3-030-01258-8_39}

\bibitem{doi:10.1109/wacvw54805.2022.00036}
\newblock \emph{Powerful Physical Adversarial Examples Against Practical Face Recognition Systems}.
\newblock DOI:10.1109/wacvw54805.2022.00036, 2022.
\newblock \url{https://doi.org/10.1109/wacvw54805.2022.00036}

\bibitem{doi:10.32604/jai.2021.014175}
\newblock \emph{An Adversarial Attack System for Face Recognition}.
\newblock DOI:10.32604/jai.2021.014175, 2021.
\newblock \url{https://doi.org/10.32604/jai.2021.014175}

\bibitem{doi:10.1109/access.2020.2993304}
\newblock \emph{Towards Adversarial Robustness via Feature Matching}.
\newblock DOI:10.1109/access.2020.2993304, 2020.
\newblock \url{https://doi.org/10.1109/access.2020.2993304}

\bibitem{doi:10.1109/cvpr46437.2021.00846}
\newblock \emph{Dual Attention Suppression Attack: Generate Adversarial Camouflage in Physical World}.
\newblock DOI:10.1109/cvpr46437.2021.00846, 2021.
\newblock \url{https://doi.org/10.1109/cvpr46437.2021.00846}

\bibitem{doi:10.1109/lra.2022.3189783}
\newblock \emph{Physical Adversarial Attack on a Robotic Arm}.
\newblock DOI:10.1109/lra.2022.3189783, 2022.
\newblock \url{https://doi.org/10.1109/lra.2022.3189783}

\bibitem{arxiv:2210.15291}
\newblock \emph{Isometric 3D Adversarial Examples in the Physical World}.
\newblock arXiv:2210.15291, 2022.
\newblock \url{https://arxiv.org/abs/2210.15291}

\bibitem{arxiv:1801.02612}
\newblock \emph{Spatially Transformed Adversarial Examples}.
\newblock arXiv:1801.02612, 2018.
\newblock \url{https://arxiv.org/abs/1801.02612}

\bibitem{arxiv:2203.08392}
\newblock \emph{Patch-Fool: Are Vision Transformers Always Robust Against Adversarial Perturbations?}.
\newblock arXiv:2203.08392, 2022.
\newblock \url{https://arxiv.org/abs/2203.08392}

\end{thebibliography}

\end{document}
